\documentclass[11pt,a4paper]{article}

% ------------------------------------------------------------------ packages
\usepackage[a4paper,left=2.6cm,right=2.6cm,top=2.7cm,bottom=2.7cm]{geometry}
\usepackage{amsmath,amssymb,amsthm,mathtools,bm}
\usepackage{booktabs,array,comment}
\usepackage[shortlabels]{enumitem}
\usepackage{microtype}
\usepackage{xcolor}
\definecolor{linkblue}{RGB}{0,0,140}
\usepackage{fancyhdr}
\usepackage[colorlinks=true,linkcolor=linkblue,citecolor=linkblue,urlcolor=linkblue]{hyperref}
\usepackage[capitalise,noabbrev]{cleveref}
\usepackage{booktabs}
\usepackage{multirow}
\usepackage{siunitx}

\numberwithin{equation}{section}
\setlength{\parindent}{1.2em}
\setlength{\parskip}{3pt plus 1pt}
\renewcommand{\arraystretch}{1.25}
\setlist{itemsep=3pt,topsep=4pt}
\allowdisplaybreaks

% ------------------------------------------------------------------ headers
\pagestyle{fancy}
\fancyhf{}
% \fancyhead[C]{\small\itshape Time-fractional Oseen equations}
\fancyfoot[C]{\thepage}
\renewcommand{\headrulewidth}{0.3pt}

% ------------------------------------------------------------------ theorems
\theoremstyle{plain}
\newtheorem{theorem}{Theorem}[section]
\newtheorem{lemma}[theorem]{Lemma}
\newtheorem{proposition}[theorem]{Proposition}
\newtheorem{corollary}[theorem]{Corollary}
\theoremstyle{definition}
\newtheorem{assumption}[theorem]{Assumption}
\theoremstyle{remark}
\newtheorem{remark}[theorem]{Remark}
\crefname{assumption}{Assumption}{Assumptions}

% ------------------------------------------------------------------ macros
\newcommand{\CD}{{}^{C}\!D_{0,t}^{\alpha}}   % Caputo derivative
\newcommand{\RL}{\partial_t^{-\alpha}}    % Riemann--Liouville integral
\newcommand{\nrm}[1]{|||#1|||_h}      % energy norm
\newcommand{\Dt}{D_\tau^{\alpha}}     % discrete Caputo derivative
\newcommand{\bmb}{\boldsymbol b}
\DeclareMathOperator{\diam}{diam}

\begin{document}

% ------------------------------------------------------------------ title
\begin{center}
{\LARGE\bfseries Symmetric stabilized finite element methods\\[3pt] for time-fractional Oseen equations\par}
\vspace{1.2em}
{\large Naveed Ahmed$^{1}$ \quad and \quad Samir Karaa$^{2}$}\par
\vspace{0.7em}
{\small
$^{1}$\, Department of Mathematics and Natural Sciences, Gulf University for Science and Technology, Kuwait\\
\texttt{ahmed.n@gust.edu.kw}\\[2pt]
$^{2}$\, Department of Mathematics, Sultan Qaboos University, Al-Khod 123, Muscat, Oman\\
\texttt{skaraa@squ.edu.om}\par}
\end{center}
\vspace{0.6em}

\begin{abstract}
\noindent
We study inf--sup stable finite element discretizations of the 
time-fractional Oseen equations, stabilized by symmetric convection 
stabilization (local projection, continuous interior penalty, orthogonal 
subscale, or subgrid scale) and grad-div stabilization, for homogeneous 
Dirichlet boundary conditions and with an optional reaction term. We prove 
stability and error estimates for the semi-discrete problem and for the L1 
time discretization. The error analysis compares the discrete solution with 
a Fortin-type interpolant rather than a Stokes/Ritz projection; the $L^2$ 
norm of the error is controlled by a fractional Gronwall argument, and 
the pressure--divergence coupling by the grad-div term. 

\end{abstract}

\noindent\textbf{Keywords:} time-fractional Oseen equations; Caputo derivative; symmetric stabilisation; grad-div stabilisation; L1 scheme.\\
\textbf{MSC 2020:} 65M60, 65M12, 35R11, 76D07.

% =====================================================================
\section{Introduction}\label{sec:intro}
We study the time-fractional Oseen equations and seek a velocity $\bm u=(u_1,u_2)^\top$ and pressure $p$ satisfying, for $(x,t)\in\Omega\times(0,T]$,
\begin{equation}\label{eq:pde}
\begin{aligned}
\CD\bm u-\nu\Delta\bm u+(\bmb\cdot\nabla)\bm u+\mu\bm u+\nabla p&=\bm f, \qquad 0<\alpha<1,\quad \nu>0,\quad \mu\ge0,\
\nabla!\cdot\bm u&=0,\
\bm u(\cdot,0)&=\bm u_0,
\end{aligned}
\end{equation}
where $\Omega$ is a bounded polygon, and the Caputo derivative of order $\alpha$ is defined by
\begin{equation}\label{eq:caputo}
\CD\bm u(x,t)=\frac{1}{\Gamma(1-\alpha)}\int_0^t(t-s)^{-\alpha},\partial_s\bm u(x,s),ds .
\end{equation}
We impose homogeneous Dirichlet conditions on the velocity, $\bm u=0$ on $\partial\Omega$, and normalise the pressure by $\int_\Omega p,dx=0$. When $\mu=0$, the \eqref{eq:pde} is the classical Oseen problem with a Caputo time derivative.

Two features shape the numerical analysis. When $\nu$ is small, convection can dominate diffusion, so stabilisation is needed to obtain estimates that stay robust with respect to viscosity. The Caputo derivative is nonlocal, and the solution may have limited temporal regularity near $t=0$; in particular, its time derivative can behave like $t^{\alpha-1}$. For this reason, standard arguments that require $\partial_t\bm u\in L^2$ may not apply. We use inf–sup stable finite element pairs with grad-div stabilisation and one of four symmetric convection stabilisations: local projection (LPS), continuous interior penalty (CIP), orthogonal subscale (OSS), or subgrid scale (SGS). Time is discretised with the L1 scheme.

For the error analysis, we compare the discrete velocity with a divergence-preserving Fortin interpolant rather than a Stokes projection. A fractional energy argument controls the $L^2$ error, while grad-div stabilisation controls the pressure–divergence coupling. Viscosity-robust estimates of this kind are available for integer-order problems, including CIP \cite{BF07} and grad-div stabilisation \cite{dFGJN}. Here we extend that approach to the Caputo setting under regularity assumptions that accommodate the initial-time singularity.

\paragraph{Notation.}
The symbols $|\cdot|$ and $(\cdot,\cdot)$ denote the $L^2(\Omega)$ norm and inner product, respectively, applied componentwise to vector fields. The advection field satisfies $\bmb\in[W^{1,\infty}(\Omega)]^2$ and is divergence-free:
\[
\nabla!\cdot\bmb=0 .
\]
The Friedrichs constant $C_F$ satisfies
\[
|\bm v|\le C_F|\nabla\bm v|\qquad\text{for all }\bm v\in[H^1_0(\Omega)]^2 .
\]
We write the Riemann–Liouville integral as
\[
\RL g(t)=\frac1{\Gamma(\alpha)}\int_0^t(t-s)^{\alpha-1}g(s),ds .
\]
Throughout, $C$ denotes a generic constant with no negative powers of $\nu$; any other dependencies are stated where relevant.

\paragraph{Outline.}
\Cref{sec:space} develops the spatial discretisation and its stabilisations, then analyses the semi-discrete problem. \Cref{sec:full} presents the L1 scheme, and the final section reports the numerical results.
% =====================================================================
\section{Space discretisation and symmetric stabilisation}\label{sec:space}
Let $\{\mathcal T_h\}$ be a shape-regular family of meshes of $\Omega$. We use triangular elements, or quadrilateral elements with $\mathbb P_r$ replaced by $\mathbb Q_r$. For each element $K$, write $h_K=\diam K$, and set
\[
h=\max_{K\in\mathcal T_h}h_K\le1 .
\]
The velocity and pressure spaces are
\[
\bm V_h=\bigl\{\bm v_h\in[H^1_0(\Omega)]^2:\ \bm v_h|_K\in[\mathbb P_r(K)]^2\ \ \forall K\in\mathcal T_h\bigr\},
\]
\[
Q_h=\bigl\{q_h\in L^2_0(\Omega):\ q_h|_K\in\mathbb P_s(K)\ \ \forall K\in\mathcal T_h\bigr\},
\]
Here $r\ge1$, and $Q_h$ may be continuous or discontinuous, depending on the finite element pair. We define the pressure--velocity coupling by
\[
b(\bm v,q)=-(\nabla\!\cdot\bm v,q).
\]
We assume that $(\bm V_h,Q_h)$ is inf--sup stable in the $H^1$ seminorm: for some $\beta_1>0$ independent of $h$, $\nu$, and $\tau$,
\begin{equation}\label{eq:infsup}
\inf_{0\ne q_h\in Q_h}\ \sup_{0\ne\bm v_h\in\bm V_h}\ \frac{b(\bm v_h,q_h)}{\|q_h\|\,\|\nabla\bm v_h\|}\ \ge\ \beta_1 .
\end{equation}

Following~\cite{AKJ}, we consider four symmetric stabilisations.
\begin{itemize}
\item \emph{Local projection (LPS).} Let $D(K)$ be a finite-dimensional local space, let $\pi_K:L^2(K)\to D(K)$ be the $L^2$ projection, and set $\kappa_K=I-\pi_K$. The stabilisation is
\[
S_h^{\rm LPS}(\bm w,\bm v)=\sum_{K\in\mathcal T_h}\tau_K\,(\kappa_K\nabla\bm w,\kappa_K\nabla\bm v)_K,\qquad \tau_K\sim h_K .
\]
One may replace the full gradient by the streamline derivative $(\bmb_K\cdot\nabla)$, with $\bmb_K$ a piecewise constant approximation of $\bmb$.
\item \emph{Continuous interior penalty (CIP).} The sum is over interior faces:
\[
S_h^{\rm CIP}(\bm w,\bm v)=\gamma_{\rm cip}\sum_{F\in\mathcal F_h^{\rm int}}\int_Fh_F^2\,|\bmb\cdot\bm n_F|\,[\![\bm n_F\!\cdot\!\nabla\bm w]\!]\cdot[\![\bm n_F\!\cdot\!\nabla\bm v]\!]\,ds .
\]
\item \emph{Orthogonal subscale (OSS).} Set $\bm e_b=\bmb/\|\bmb\|_{\ell^2}$ and define the weighted $L^2$ projection $\pi_h^k$ by
\[
\bigl(h\|\bmb\|_{\ell^2}\,\pi_h^kg,\bm\xi_h\bigr)=\bigl(h\|\bmb\|_{\ell^2}\,g,\bm\xi_h\bigr)\qquad\text{for all projection test functions }\bm\xi_h .
\]
The corresponding stabilisation is
\[
S_h^{\rm OSS}(\bm w,\bm v)=\tau_{\rm oss}\Bigl(h\|\bmb\|_{\ell^2}\,(I-\pi_h^k)(\bm e_b\!\cdot\!\nabla\bm w),\,(I-\pi_h^k)(\bm e_b\!\cdot\!\nabla\bm v)\Bigr),
\]
This form is symmetric and nonnegative.
\item \emph{Subgrid scale (SGS).} Given a scale-separation projection $\Pi$, define
\[
S_h^{\rm SGS}(\bm w,\bm v)=\sum_{K\in\mathcal T_h}\tau_K\int_K\nabla\Pi(\bm w):\nabla\Pi(\bm v)\,dx .
\]
\end{itemize}

\begin{assumption}[Fortin-type interpolant]\label{ass:I}
We assume there is a linear, time-independent operator $\bm I_h:[H^{r+1}\cap H^1_0]^2\to\bm V_h$, bounded on $[H^1_0]^2$, that satisfies the following properties for every admissible $\bm v$:
\begin{enumerate}[(F1),leftmargin=2.8em]
\item Fortin property:
\[
b(\bm v-\bm I_h\bm v,q_h)=0\qquad\text{for all }q_h\in Q_h ;
\]
\item gradient approximation:
\[
\|\nabla(\bm v-\bm I_h\bm v)\|\le C_Ih^r\|\bm v\|_{r+1} ;
\]
\item $L^2$ approximation:
\[
\|\bm v-\bm I_h\bm v\|\le C_Ih^{r+1}\|\bm v\|_{r+1} ;
\]
\item pressure approximation, for $q\in H^{s+1}\cap L^2_0$:
\[
\inf_{q_h\in Q_h}\|q-q_h\|\le C_Ih^{s+1}\|q\|_{s+1} .
\]
\end{enumerate}
\end{assumption}
For standard finite element pairs, Fortin's criterion gives (F1) and (F2) from \eqref{eq:infsup}; see \cite{BBF}. Property (F3) is an additional approximation assumption. It holds, for example, for the quasi-local, divergence-preserving interpolant of Girault and Scott \cite{GS03} for the Taylor--Hood family. We will also use
\[
\|\bm w\|\le C_F\|\nabla\bm w\|\le C_FC_Ih^r\|\bm u\|_{r+1}
\]
This bound is independent of negative powers of $\nu$, but it gives one order less in the lower-order terms; see \cref{rem:hyp}(ii). Since $\bm I_h$ is linear, time independent, and bounded on $H^1$, it commutes with the Caputo derivative:
\[
\CD(\bm I_h\bm u)=\bm I_h\CD\bm u .
\]
We apply $\bm I_h$ only to the velocity.

\begin{assumption}[Stabilisation]\label{ass:S}
The stabilisation $S_h$ is a symmetric, positive semidefinite bilinear form on $\bm V_h+[H^{r+1}\cap H^1_0]^2$. With the interpolant $\bm I_h$ from \cref{ass:I}, we assume:
\begin{enumerate}[(S1),leftmargin=2.8em]
\item for all $\bm v_h\in\bm V_h$:
\[
S_h(\bm v_h,\bm v_h)\le C_Sh\,\|\nabla\bm v_h\|^2 ;
\]
\item for all $\bm u\in[H^{r+1}\cap H^1_0]^2$:
\[
S_h(\bm u,\bm u)+S_h(\bm I_h\bm u,\bm I_h\bm u)\le C_Sh^{2r+1}\|\bm u\|_{r+1}^2 .
\]
\end{enumerate}
\end{assumption}
The standard LPS, CIP, OSS, and SGS choices satisfy these assumptions when the associated projections and interpolants have the stated approximation properties.
\[
S_h(\bm u,\bm u)\lesssim\sum_K\tau_Kh_K^{2r}|\bm u|_{r+1,K}^2 .
\]
For CIP, the jump $[\![\bm n_F\!\cdot\!\nabla\bm u]\!]$ vanishes when $\bm u\in H^2$. A trace estimate applied to $\bm I_h\bm u-\bm u$ then gives (S2). This step uses the local bound $h_K\|\nabla^2(\bm I_h\bm u-\bm u)\|_K\le Ch^r\|\bm u\|_{r+1,\omega_K}$, which holds for quasi-local interpolants such as the one in \cite{GS03}. The OSS and SGS cases follow by similar arguments. The constants depend on $\|\bmb\|_{L^\infty}$, $\gamma_{\rm cip}$, and $\tau_{\rm oss}$.

The semi-discrete problem is to find $(\bm u_h(t),p_h(t))\in\bm V_h\times Q_h$, with initial value $\bm u_h(0)=\bm u_{h,0}$, such that for every $(\bm v_h,q_h)\in\bm V_h\times Q_h$,
\begin{equation}\label{eq:semi}
\begin{aligned}
\bigl(\CD\bm u_h,\bm v_h\bigr)+a_h(\bm u_h,\bm v_h)+b(\bm v_h,p_h)&=(\bm f,\bm v_h),\\
b(\bm u_h,q_h)&=0 
\end{aligned}
\end{equation}
Here the unstabilized bilinear form is
\begin{equation}\label{eq:forms}
a(\bm w,\bm v)=\nu(\nabla\bm w,\nabla\bm v)+\bigl((\bmb\cdot\nabla)\bm w,\bm v\bigr)+\mu(\bm w,\bm v).
\end{equation}    
We add the grad-div and convection-stabilisation terms to obtain
\begin{align*}
a_h(\bm w,\bm v)&=a(\bm w,\bm v)+\gamma(\nabla\!\cdot\bm w,\nabla\!\cdot\bm v)+S_h(\bm w,\bm v)
\end{align*}
We use the energy norm
\[
\nrm{\bm v}^2=\nu\|\nabla\bm v\|^2+\mu\|\bm v\|^2+\gamma\|\nabla\!\cdot\bm v\|^2+S_h(\bm v,\bm v).
\]
To see that \eqref{eq:semi} is well posed, restrict the problem to the discretely divergence-free subspace $\{\bm v_h\in\bm V_h:\ b(\bm v_h,q_h)=0\ \forall q_h\in Q_h\}$. On this subspace, the velocity solves a finite-dimensional fractional ODE system $M\,\CD\bm U+A\bm U=\bm F$, where the mass matrix $M$ is symmetric positive definite. Its solution is continuous on $[0,T]$ and absolutely continuous when $\bm f\in L^2(0,T;L^2)$; see \cite[Ch.\ 6]{Die10} and \cite{LL18}. The pressure is recovered from the first equation of \eqref{eq:semi} using \eqref{eq:infsup}. In particular, $\bm u_h\in W^{1,1}(0,T;\bm V_h)$, the regularity needed below.
\begin{lemma}[Coercivity]\label{lem:bounds}
\begin{enumerate}[(i),leftmargin=2.4em]
\item For $\bm w,\bm v\in[H^1_0(\Omega)]^2$,
\[
((\bmb\cdot\nabla)\bm w,\bm v)=-((\bmb\cdot\nabla)\bm v,\bm w).
\]
In particular $((\bmb\cdot\nabla)\bm v,\bm v)=0$ and $a_h(\bm v,\bm v)=\nrm{\bm v}^2$.
\item Assume $S_h(\bm v_h,\bm v_h)\le C_Sh\|\nabla\bm v_h\|^2$ on $\bm V_h$ and put
\[
M_h:=\nu^{1/2}+C_F\mu^{1/2}+\sqrt2\,\gamma^{1/2}+(C_Sh)^{1/2}.
\]
Then, for $\bm w\in\bm V_h+[H^{r+1}\cap H^1_0]^2$ (so that $S_h(\bm w,\bm w)$ is defined), $\bm v_h\in\bm V_h$ and $q\in L^2_0(\Omega)$,
\begin{align}
|b(\bm v_h,q)|&\le\sqrt2\,\|\nabla\bm v_h\|\,\|q\|, \label{eq:b-bound}\\
|a_h(\bm w,\bm v_h)|&\le\bigl(M_h\nrm{\bm w}+\|\bmb\|_{L^\infty}\|\bm w\|\bigr)\,\|\nabla\bm v_h\|, \label{eq:ah-bound}\\
\bigl|a(\bm w,\bm v_h)+\gamma(\nabla\!\cdot\bm w,\nabla\!\cdot\bm v_h)\bigr|&\le\bigl[(\nu+2\gamma+C_F\|\bmb\|_{L^\infty})\|\nabla\bm w\|+C_F\mu\|\bm w\|\bigr]\|\nabla\bm v_h\|. \label{eq:a-bound}
\end{align}

\end{enumerate}
\end{lemma}
\begin{proof}
(i) By the product rule, $\bmb\cdot\nabla(\bm w\cdot\bm v)=(\bmb\cdot\nabla\bm w)\cdot\bm v+(\bmb\cdot\nabla\bm v)\cdot\bm w$. Integrating and applying the divergence theorem,
\[
((\bmb\cdot\nabla)\bm w,\bm v)+((\bmb\cdot\nabla)\bm v,\bm w)=\int_\Omega\bmb\cdot\nabla(\bm w\cdot\bm v)\,dx
=-\int_\Omega(\nabla\!\cdot\bmb)\,\bm w\cdot\bm v\,dx+\int_{\partial\Omega}(\bmb\cdot\bm n)\,\bm w\cdot\bm v\,ds .
\]
The first term on the right vanishes because $\nabla\!\cdot\bmb=0$. The boundary term vanishes because $\bm w=\bm v=0$ on $\partial\Omega$; no condition on $\bmb\cdot\bm n$ is needed. Taking $\bm w=\bm v$ gives $2((\bmb\cdot\nabla)\bm v,\bm v)=0$. Consequently, from the definitions in \eqref{eq:forms},
\[
a_h(\bm v,\bm v)=\nu\|\nabla\bm v\|^2+\underbrace{((\bmb\cdot\nabla)\bm v,\bm v)}_{=0}+\mu\|\bm v\|^2+\gamma\|\nabla\!\cdot\bm v\|^2+S_h(\bm v,\bm v)=\nrm{\bm v}^2 .
\]

(ii) Pointwise, $|\nabla\!\cdot\bm v|^2=(\partial_1v_1+\partial_2v_2)^2\le2\bigl((\partial_1v_1)^2+(\partial_2v_2)^2\bigr)\le2|\nabla\bm v|^2$, hence $\|\nabla\!\cdot\bm v\|\le\sqrt2\|\nabla\bm v\|$ and
\[
|b(\bm v_h,q)|=|(\nabla\!\cdot\bm v_h,q)|\le\|\nabla\!\cdot\bm v_h\|\,\|q\|\le\sqrt2\,\|\nabla\bm v_h\|\,\|q\|,
\]
which is \eqref{eq:b-bound}. For \eqref{eq:ah-bound} we bound the terms of $a_h(\bm w,\bm v_h)$ one by one, using the definition of $\nrm{\cdot}$ and the Cauchy--Schwarz inequality:
\[
\begin{aligned}
\nu(\nabla\bm w,\nabla\bm v_h)&\le\nu^{1/2}\bigl(\nu^{1/2}\|\nabla\bm w\|\bigr)\|\nabla\bm v_h\|\le\nu^{1/2}\nrm{\bm w}\,\|\nabla\bm v_h\|,\\
\mu(\bm w,\bm v_h)&\le\mu^{1/2}\bigl(\mu^{1/2}\|\bm w\|\bigr)\|\bm v_h\|\le\mu^{1/2}\nrm{\bm w}\,C_F\|\nabla\bm v_h\|,\\
\gamma(\nabla\!\cdot\bm w,\nabla\!\cdot\bm v_h)&\le\gamma^{1/2}\bigl(\gamma^{1/2}\|\nabla\!\cdot\bm w\|\bigr)\|\nabla\!\cdot\bm v_h\|\le\sqrt2\,\gamma^{1/2}\nrm{\bm w}\,\|\nabla\bm v_h\|,\\
S_h(\bm w,\bm v_h)&\le S_h(\bm w,\bm w)^{1/2}S_h(\bm v_h,\bm v_h)^{1/2}\le(C_Sh)^{1/2}\nrm{\bm w}\,\|\nabla\bm v_h\|.
\end{aligned}
\]
For the convection term, we use (i), moving the derivative onto $\bm v_h$:
\[
((\bmb\cdot\nabla)\bm w,\bm v_h)=-((\bmb\cdot\nabla)\bm v_h,\bm w)\le\|\bmb\|_{L^\infty}\|\nabla\bm v_h\|\,\|\bm w\|.
\]
Adding the five bounds gives
\[
|a_h(\bm w,\bm v_h)|\le\Bigl[\nu^{1/2}+C_F\mu^{1/2}+\sqrt2\gamma^{1/2}+(C_Sh)^{1/2}\Bigr]\nrm{\bm w}\,\|\nabla\bm v_h\|+\|\bmb\|_{L^\infty}\|\bm w\|\,\|\nabla\bm v_h\|,
\]
which is \eqref{eq:ah-bound}. For \eqref{eq:a-bound} we estimate directly
\[
\begin{aligned}
\nu(\nabla\bm w,\nabla\bm v_h)&\le\nu\|\nabla\bm w\|\,\|\nabla\bm v_h\|,\\
((\bmb\cdot\nabla)\bm w,\bm v_h)&\le\|\bmb\|_{L^\infty}\|\nabla\bm w\|\,\|\bm v_h\|\le C_F\|\bmb\|_{L^\infty}\|\nabla\bm w\|\,\|\nabla\bm v_h\|,\\
\mu(\bm w,\bm v_h)&\le\mu\|\bm w\|\,\|\bm v_h\|\le C_F\mu\|\bm w\|\,\|\nabla\bm v_h\|,\\
\gamma(\nabla\!\cdot\bm w,\nabla\!\cdot\bm v_h)&\le\gamma\|\nabla\!\cdot\bm w\|\,\|\nabla\!\cdot\bm v_h\|\le2\gamma\|\nabla\bm w\|\,\|\nabla\bm v_h\|,
\end{aligned}
\]
and add.
\end{proof}

% =====================================================================
% \subsection{Fractional calculus preliminaries}\label{sec:tools}
For absolutely continuous $y$ and nonnegative, integrable $g$, we will use the following identities and inequalities:
\begin{equation}\label{eq:rl}
\begin{aligned}
&\RL\CD y=y-y(0),\qquad \RL1=\frac{t^\alpha}{\Gamma(\alpha+1)},\qquad \RL g\ge0,\\
&\RL\RL g\le\frac{t^\alpha\,\Gamma(\alpha)}{\Gamma(2\alpha)}\,\RL g,\qquad(\RL g)^2\le\frac{t^\alpha}{\Gamma(\alpha+1)}\,\RL g^2 .
\end{aligned}
\end{equation}
The fourth relation follows from 
$$\RL\RL g=\frac1{\Gamma(2\alpha)}\int_0^t(t-s)^{2\alpha-1}g(s)\,ds$$ 
and $$(t-s)^{2\alpha-1}=(t-s)^\alpha(t-s)^{\alpha-1}\le t^\alpha(t-s)^{\alpha-1}$$
The fifth inequality is Cauchy--Schwarz with respect to the measure $(t-s)^{\alpha-1}ds/\Gamma(\alpha)$, whose total mass is $\RL1$. For vector-valued $\bm g$, the triangle inequality gives $\|\RL\bm g\|\le\RL\|\bm g\|$; likewise, $\nrm{\RL\bm g}\le\RL\nrm{\bm g}$ for the seminorm $\nrm{\cdot}$. The first three identities are standard; see \cite[Ch.\ 2--3]{Die10}.

\begin{lemma}[Positivity of the Caputo derivative \cite{Ali}]\label{lem:posc}
If $\bm v:[0,T]\to L^2(\Omega)^2$ is absolutely continuous, then
\[
(\CD\bm v,\bm v)\ge\tfrac12\CD\|\bm v\|^2 .
\]
\end{lemma}

\begin{lemma}[Fractional Gronwall inequality; see \cite{YGD07,LL18}]\label{lem:gronwall}
Let $y\ge0$ be absolutely continuous, $g\ge0$ integrable, $c_1>0$, and
\[
\CD y\le c_1y+g .
\]
Then
\[
y(t)\le y(0)\,E_\alpha(c_1t^\alpha)+\Gamma(\alpha)\,E_{\alpha,\alpha}(c_1t^\alpha)\,\RL g(t).
\]
Consequently, with
\[
d:=\max\bigl\{E_\alpha(T^\alpha),\ \Gamma(\alpha)E_{\alpha,\alpha}(T^\alpha)\bigr\}\ge1
\]
and $c_1=1$,
\[
y(t)\le d\bigl(y(0)+\RL g(t)\bigr),\qquad
\RL y(t)\le t^\alpha d\Bigl(\frac{y(0)}{\Gamma(\alpha+1)}+\frac{\Gamma(\alpha)}{\Gamma(2\alpha)}\RL g(t)\Bigr).
\]
\end{lemma}

We will use the fractional Gronwall inequality in \cref{sec:full}, where its discrete counterpart is needed. For the semi-discrete analysis, the following simpler estimate suffices. It avoids Mittag--Leffler constants and handles the forcing term $a$ through $\RL a$, rather than $\RL a^2$.
\begin{lemma}[Energy estimate without Gronwall]\label{lem:sup}
Let $y\in C([0,T])$, $y\ge0$, with $y^2$ absolutely continuous, let $z,a,g\ge0$ be integrable, and suppose that
\[
\CD(y^2)+z\le 2ay+g\qquad\text{a.e.\ in }(0,T).
\]
Set $\bar a(t):=\sup_{0<s\le t}\RL a(s)$ and $\bar g(t):=\sup_{0<s\le t}\RL g(s)$, assumed finite. Then, for $t\in(0,T]$,
\begin{equation}\label{eq:sup}
\sup_{0\le s\le t}y(s)\le y(0)+2\bar a(t)+\bar g(t)^{1/2},\qquad
\RL z(t)\le y(0)^2+\bar g(t)+2\bar a(t)\sup_{0\le s\le t}y(s).
\end{equation}
\end{lemma}
\begin{proof}
Let $Y(t):=\sup_{0\le s\le t}y(s)$. Applying $\RL$ and using $\RL\CD(y^2)=y^2-y(0)^2$ from \eqref{eq:rl} gives
\[
y(t)^2+\RL z(t)\le y(0)^2+2\,\RL(ay)(t)+\RL g(t)\le y(0)^2+2Y(t)\,\bar a(t)+\bar g(t),
\]
since the kernel of $\RL$ is positive and $y\le Y(t)$ on $[0,t]$. The right-hand side is nondecreasing in $t$, so the same bound holds for $y(s)^2$ for every $s\le t$, and taking the supremum over $s$ yields $Y^2\le y(0)^2+\bar g+2\bar aY$. Solving this quadratic inequality,
\[
Y\le\bar a+\sqrt{\bar a^2+y(0)^2+\bar g}\le y(0)+2\bar a+\bar g^{1/2},
\]
which is the first claim; the second is the displayed inequality with $y(t)^2\ge0$ dropped.
\end{proof}



\subsection{Stability analysis}
We first estimate the semi-discrete velocity and pressure. Testing with the velocity and using coercivity gives the energy bound; the pressure estimate then follows from the momentum equation and the discrete inf--sup condition.

% =====================================================================
% \subsection{Stability of the semi-discrete problem}\label{sec:stab}
\begin{theorem}[Semi-discrete stability]\label{thm:stab}
Let $(\bm u_h,p_h)$ solve \eqref{eq:semi} with $\bm f\in L^2(0,T;L^2)$. Then, for almost every $t\in(0,T]$,
\begin{equation}\label{eq:energy}
\tfrac12\CD\|\bm u_h\|^2+\nrm{\bm u_h}^2\le(\bm f,\bm u_h),
\end{equation}
and, if $F(t):=\sup_{0<s\le t}\RL\|\bm f\|(s)<\infty$ (for example, if $\bm f\in L^p(0,T;L^2)$ with $p>1/\alpha$),
\begin{equation}\label{eq:stab}
\sup_{0\le s\le t}\|\bm u_h(s)\|\le\|\bm u_{h,0}\|+2F(t),\qquad
2\,\RL\nrm{\bm u_h}^2(t)\le\|\bm u_{h,0}\|^2+2F(t)\bigl(\|\bm u_{h,0}\|+2F(t)\bigr).
\end{equation}
For $\bm f\in L^2(0,T;L^2)$ alone, \cref{lem:gronwall} instead gives the pointwise estimate
$\|\bm u_h(t)\|^2\le\|\bm u_{h,0}\|^2E_\alpha(t^\alpha)+\Gamma(\alpha)E_{\alpha,\alpha}(t^\alpha)\,\RL\|\bm f\|^2(t)$.
If $\mu>0$, the Gronwall lemma can be avoided:
\[
\|\bm u_h(t)\|^2+\RL\nrm{\bm u_h}^2\le\|\bm u_{h,0}\|^2+\mu^{-1}\RL\|\bm f\|^2 .
\]
\end{theorem}
\begin{proof}
Test \eqref{eq:semi} with $\bm v_h=\bm u_h$ and $q_h=p_h$. Since the discrete divergence constraint gives $b(\bm u_h,p_h)=0$, the first equation becomes
\[
(\CD\bm u_h,\bm u_h)+a_h(\bm u_h,\bm u_h)=(\bm f,\bm u_h).
\]
By \cref{lem:posc} and \cref{lem:bounds}(i), $(\CD\bm u_h,\bm u_h)\ge\frac12\CD\|\bm u_h\|^2$ and $a_h(\bm u_h,\bm u_h)=\nrm{\bm u_h}^2$; this is \eqref{eq:energy}.

Multiplying \eqref{eq:energy} by $2$ and applying Cauchy--Schwarz gives
\[
\CD\|\bm u_h\|^2+2\nrm{\bm u_h}^2\le2\|\bm f\|\,\|\bm u_h\| ,
\]
This is the hypothesis of \cref{lem:sup} with $y=\|\bm u_h\|$, $z=2\nrm{\bm u_h}^2$, $a=\|\bm f\|$, and $g=0$, so \eqref{eq:sup} yields \eqref{eq:stab}. For the pointwise estimate, use $2(\bm f,\bm u_h)\le\|\bm f\|^2+\|\bm u_h\|^2$, discard the nonnegative energy term, and apply \cref{lem:gronwall} with $y=\|\bm u_h\|^2$, $c_1=1$, and $g=\|\bm f\|^2$.

If $\mu>0$, then $\|\bm u_h\|\le\mu^{-1/2}\nrm{\bm u_h}$, so that
\[
(\bm f,\bm u_h)\le\|\bm f\|\,\|\bm u_h\|\le\mu^{-1/2}\|\bm f\|\,\nrm{\bm u_h}\le\tfrac12\nrm{\bm u_h}^2+\tfrac1{2\mu}\|\bm f\|^2 .
\]
Inserting this into \eqref{eq:energy} and multiplying by $2$ gives $\CD\|\bm u_h\|^2+\nrm{\bm u_h}^2\le\mu^{-1}\|\bm f\|^2$; applying $\RL$ proves the last claim.
\end{proof}
\begin{remark}
Estimate \eqref{eq:stab} contains neither exponential growth in $T$ nor negative powers of $\nu$. The advection field enters only through the skew-symmetry in \cref{lem:bounds}(i). A sharper estimate, $\|\bm u_h(t)\|\le\|\bm u_{h,0}\|+\RL\|\bm f\|(t)$, also holds for $\bm f\in L^1(0,T;L^2)$: the discrete Oseen operator is accretive on the discretely divergence-free subspace, and subordination makes the associated solution operators contractions; see \cite{Baz01}. We do not use this sharper bound below.
\end{remark}

% ---- pressure stability     ppp
%\newpage

The velocity energy estimate does not directly control the discrete pressure. We therefore combine the momentum equation with the discrete inf--sup condition to obtain the following bound for $p_h$.
%
\begin{theorem}[Pressure stability]
\label{thm:semidiscrete_pressure_stability}
Let $(\bm u_h,p_h)$ solve \eqref{eq:semi} with 
$\bm f \in L^2(0,T;[L^2(\Omega)]^2).$
Then, for a.e.\ $t\in(0,T]$,
\[
\begin{aligned}
\beta_1\|p_h(t)\|
\le {}&
C_F\|\bm f(t)\|
+C_F\|{}^{C}D_{0,t}^{\alpha}\bm u_h(t)\|
\\
&\quad
+M_h|||\bm u_h(t)|||_h
+\|\boldsymbol b\|_{L^\infty}\|\bm u_h(t)\|.
\end{aligned}
\]
For every $t\in(0,T]$, the fractionally integrated pressure satisfies the following bound:
\begin{equation}\label{eq:pstab-int}
\beta_1\|\RL p_h(t)\|\le C_F\bigl(\|\bm u_h(t)\|+\|\bm u_{h,0}\|\bigr)+C_F\,\RL\|\bm f\|(t)+M_h\,\RL\nrm{\bm u_h}(t)+\|\bmb\|_{L^\infty}\RL\|\bm u_h\|(t),
\end{equation}
The terms on the right are controlled by \cref{thm:stab} and \eqref{eq:rl}.
\end{theorem}
%
\begin{proof}
By the discrete inf--sup condition,
\[
\beta_1\|p_h(t)\|
\le
\sup_{0\ne \bm v_h\in \bm  V_h}
\frac{|b(\bm v_h,p_h(t))|}{\|\nabla \bm v_h\|}.
\]
From the first equation of the semidiscrete problem,
\[
b(\bm v_h,p_h)
=
(\bm f,\bm v_h)
-\bigl({}^{C}D_{0,t}^{\alpha}\bm u_h,\bm v_h\bigr)
-a_h(\bm u_h,\bm v_h).
\]
Therefore,
\[
\begin{aligned}
|b(\bm v_h,p_h)|
\le {}&
|(\bm f,\bm v_h)|
+
\bigl|({}^{C}D_{0,t}^{\alpha}\bm u_h,\bm v_h)\bigr|
+
|a_h(\bm u_h,\bm v_h)|.
\end{aligned}
\]
By the Friedrichs inequality,
\[
|(\bm f,\bm v_h)|
\le
C_F\|\bm f\|\,\|\nabla \bm v_h\|,
\]
and
\[
\bigl|({}^{C}D_{0,t}^{\alpha}\bm u_h,\bm v_h)\bigr|
\le
C_F\|{}^{C}D_{0,t}^{\alpha}\bm u_h\|
\,\|\nabla \bm v_h\|.
\]
Furthermore, the continuity estimate for $a_h$ gives
\[
|a_h(\bm u_h,\bm v_h)|
\le
\left(
M_h|||\bm u_h|||_h
+
\|\boldsymbol b\|_{L^\infty}\|\bm u_h\|
\right)
\|\nabla \bm v_h\|.
\]
Combining these estimates yields
\[
\begin{aligned}
|b(\bm v_h,p_h)|
\le
\Bigl(
&C_F\|\bm f\|
+C_F\|{}^{C}D_{0,t}^{\alpha}\bm u_h\|
\\
&+
M_h|||\bm u_h|||_h
+\|\boldsymbol b\|_{L^\infty}\|\bm u_h\|
\Bigr)
\|\nabla \bm v_h\|.
\end{aligned}
\]
Taking the supremum over $0\ne \bm v_h\in \bm V_h$ and applying the discrete inf--sup condition proves the pointwise estimate. To obtain \eqref{eq:pstab-int}, fix $\bm v_h$ and apply $\RL$ to the identity for $b(\bm v_h,p_h)$. By \eqref{eq:rl}, $\RL\CD\bm u_h=\bm u_h-\bm u_{h,0}$. Bound $a_h(\RL\bm u_h,\bm v_h)$ using \eqref{eq:ah-bound}, together with $\nrm{\RL\bm u_h}\le\RL\nrm{\bm u_h}$ and $\|\RL\bm u_h\|\le\RL\|\bm u_h\|$. Since $\RL p_h(t)\in Q_h$, another application of the inf--sup condition gives \eqref{eq:pstab-int}.
\end{proof}
The pointwise pressure estimate is conditional because it contains $\|\CD\bm u_h\|$, which is not controlled by the energy bound without introducing a negative power of $\nu$. Indeed, testing \eqref{eq:semi} with $\bm v_h=\CD\bm u_h$ yields a bound for $\RL\|\CD\bm u_h\|^2$ that involves $\nu^{-1}\|\bmb\|_{L^\infty}^2$. In contrast, \eqref{eq:pstab-int} gives a viscosity-robust estimate for the fractionally integrated pressure, analogous to the integer-order estimate in \cite{BF07}.




%=====================================================================

\subsection{Error analysis}

%=====================================================================

% =====================================================================
Let $(\bm u,p)$ be the exact solution. We assume
\begin{align}\label{sec:error}
\bm u\in & W^{1,1}\bigl(0,T;[H^1_0]^2\bigr)\cap L^\infty\bigl(0,T;[H^{r+1}\cap H^1_0]^2\bigr),\qquad
\CD\bm u\in L^\infty\bigl(0,T;[H^{r+1}\cap H^1_0]^2\bigr),\nonumber \\
&p\in L^\infty\bigl(0,T;H^{s+1}\cap L^2_0\bigr).
\end{align}
These regularity conditions depend on the domain $\Omega$ and on compatibility of the data. In time, they require only $\partial_t\bm u\in L^1$, which allows the typical initial behavior $\partial_t\bm u\sim t^{\alpha-1}$. Since $\bm u=0$ on $\partial\Omega$ for all $t$, we also have $\CD\bm u(t)\in H^1_0$. Set
\[
\bm w=\bm u-\bm I_h\bm u,\qquad \bm\zeta=\bm u_h-\bm I_h\bm u,\qquad\text{so that}\qquad \bm u_h-\bm u=\bm\zeta-\bm w ,
\]
We also let $q_h(t)\in Q_h$ be the $L^2$ projection of $p(t)$ onto $Q_h$. Then $q_h$ is measurable in time, and $\|p-q_h\|$ is the best approximation error in (F4). Constants below depend only on $C_I$ and $C_S$; $\nu$, $\gamma$, $\mu$, and $\|\bmb\|_{L^\infty}$ enter only through $c_h$ and $M_h$. Define
\begin{equation}\label{eq:vartheta}
\begin{aligned}
\vartheta_h(t)&:=c_h\,h^r\|\bm u(t)\|_{r+1}+\gamma^{-1/2}h^{s+1}\|p(t)\|_{s+1},\\
c_h&:=\|\bmb\|_{L^\infty}+\nu^{1/2}+\gamma^{1/2}+h^{1/2}+\mu h+\mu^{1/2}h,\\
\kappa_h(t)&:=h^{r+1}\|\CD\bm u(t)\|_{r+1}.
\end{aligned}
\end{equation}

\begin{lemma}[Interpolation quantities]\label{lem:interp}
Under \cref{ass:S,ass:I}, for a.e.\ $t$ and with constants depending only on $C_I$ and $C_S$:
\begin{enumerate}[(i),leftmargin=2.4em]
\item the following bounds hold:
\[
\|\CD\bm w\|\le C_I\kappa_h,\qquad \nrm{\bm w}^2\le C\vartheta_h^2,\qquad \|\bm w\|\le C\vartheta_h ;
\]
\item the quantity
\[
A:=\|\CD\bm w\|+\|\bmb\|_{L^\infty}\|\nabla\bm w\|+\mu\|\bm w\|
\]
satisfies $A\le C(\kappa_h+\vartheta_h)$;
\item the quantity
\[
B:=\nu^{1/2}\|\nabla\bm w\|+\gamma^{1/2}\|\nabla\!\cdot\bm w\|+S_h(\bm I_h\bm u,\bm I_h\bm u)^{1/2}+\gamma^{-1/2}\|p-q_h\|
\]
satisfies $B\le C\vartheta_h$.
\end{enumerate}
\end{lemma}
\begin{proof}
Note first that each summand of $c_h$ is bounded by $c_h$, and $c_h^2\ge\nu+\gamma+h+\mu h^2$ since the summands are nonnegative.

(i) Since $\bm I_h$ commutes with the Caputo derivative, $\CD\bm w=(I-\bm I_h)\CD\bm u$. Our regularity assumption ensures $\CD\bm u(t)\in[H^{r+1}\cap H^1_0]^2$, so (F3) gives $\|\CD\bm w\|\le C_Ih^{r+1}\|\CD\bm u\|_{r+1}=C_I\kappa_h$. Again by (F3) and $h\le1$,
\[
\|\bm w\|\le C_Ih^{r+1}\|\bm u\|_{r+1}\le C_Ih^{1/2}\,h^r\|\bm u\|_{r+1}\le C_Ic_h\,h^r\|\bm u\|_{r+1}\le C_I\vartheta_h .
\]
For the energy norm, we bound the four terms of $\nrm{\bm w}^2$ using (F2), (F3), $\|\nabla\!\cdot\bm w\|\le\sqrt2\|\nabla\bm w\|$ and (S2), together with 
$$S_h(\bm w,\bm w)\le2S_h(\bm u,\bm u)+2S_h(\bm I_h\bm u,\bm I_h\bm u),$$
to get
$$\nrm{\bm w}^2\le C(\nu+\mu h^2+\gamma+h)\,h^{2r}\|\bm u\|_{r+1}^2\le Cc_h^2h^{2r}\|\bm u\|_{r+1}^2\le C\vartheta_h^2.$$
(ii) The three terms of $A$ are bounded as follows:
\[
\begin{aligned}
\|\CD\bm w\|&\le C_I\kappa_h,\\
\|\bmb\|_{L^\infty}\|\nabla\bm w\|&\le C_I\|\bmb\|_{L^\infty}h^r\|\bm u\|_{r+1}\le C_Ic_h\,h^r\|\bm u\|_{r+1},\\
\mu\|\bm w\|&\le C_I\,\mu h\cdot h^r\|\bm u\|_{r+1}\le C_Ic_h\,h^r\|\bm u\|_{r+1}.
\end{aligned}
\]
Hence $A\le C_I\kappa_h+2C_Ic_hh^r\|\bm u\|_{r+1}\le C(\kappa_h+\vartheta_h)$.

(iii) The four terms of $B$ are bounded as follows:
\[
\begin{aligned}
\nu^{1/2}\|\nabla\bm w\|
+ \gamma^{1/2}\|\nabla\!\cdot\bm w\| &\le C_I \left(\nu^{1/2} + \sqrt{2} \gamma^{1/2} \right) h^r\|\bm u\|_{r+1} \\
S_h(\bm I_h\bm u,\bm I_h\bm u)^{1/2}&\le C_S^{1/2}h^{r+1/2}\|\bm u\|_{r+1}=C_S^{1/2}h^{1/2}\,h^r\|\bm u\|_{r+1},\\
\gamma^{-1/2}\|p-q_h\|&\le C_I\gamma^{-1/2}h^{s+1}\|p\|_{s+1}.
\end{aligned}
\]
Each of the first three is at most $C\,c_hh^r\|\bm u\|_{r+1}$, and the last is part of $\vartheta_h$; hence $B\le C\vartheta_h$.
\end{proof}

\begin{theorem}[Velocity error]\label{thm:err}
Under \cref{ass:S,ass:I} and \eqref{sec:error}, for $0<t\le T$,
\begin{equation}\label{eq:err-main}
\begin{aligned}
&\sup_{0\le s\le t}\|\bm u_h(s)-\bm u(s)\|+\Bigl(\RL\nrm{\bm u_h-\bm u}^2(t)\Bigr)^{1/2}\\
&\qquad\le C\Bigl[\|\bm u_{h,0}-\bm I_h\bm u_0\|+\sup_{0<s\le t}\RL(\kappa_h+\vartheta_h)(s)+\sup_{0<s\le t}\bigl(\RL\vartheta_h^2(s)\bigr)^{1/2}+\sup_{0\le s\le t}\vartheta_h(s)\Bigr],
\end{aligned}
\end{equation}
where $C$ depends only on $C_I$ and $C_S$; in particular it is independent of $h$, $\nu$, $\gamma$, $\mu$, $\alpha$ and $T$. For $\bm u_{h,0}=\bm I_h\bm u_0$ the first term vanishes; in general
\[
\|\bm u_{h,0}-\bm I_h\bm u_0\|\le\|\bm u_{h,0}-\bm u_0\|+\|\bm w(0)\| .
\]
\end{theorem}
\begin{proof}
The exact solution satisfies, for all $\bm v\in[H^1_0]^2$,
\[
(\CD\bm u,\bm v)+a(\bm u,\bm v)+b(\bm v,p)=(\bm f,\bm v).
\]
Take $\bm v=\bm v_h\in\bm V_h$ and subtract this from the first equation in \eqref{eq:semi}:
\[
\bigl(\CD(\bm u_h-\bm u),\bm v_h\bigr)+a_h(\bm u_h,\bm v_h)-a(\bm u,\bm v_h)+b(\bm v_h,p_h-p)=0 .
\]
Insert $\bm u_h-\bm u=\bm\zeta-\bm w$ and $\bm u_h=\bm\zeta+\bm I_h\bm u$:
\[
(\CD\bm\zeta,\bm v_h)-(\CD\bm w,\bm v_h)+a_h(\bm\zeta,\bm v_h)+\bigl[a_h(\bm I_h\bm u,\bm v_h)-a(\bm u,\bm v_h)\bigr]+b(\bm v_h,p_h-p)=0 .
\]
The fourth term on the left-hand side is computed from the definition of $a_h$ and $\nabla\!\cdot\bm I_h\bm u=\nabla\!\cdot\bm u-\nabla\!\cdot\bm w=-\nabla\!\cdot\bm w$:
\[
\begin{aligned}
a_h(\bm I_h\bm u,\bm v_h)-a(\bm u,\bm v_h)&=a(\bm I_h\bm u,\bm v_h)-a(\bm u,\bm v_h)+\gamma(\nabla\!\cdot\bm I_h\bm u,\nabla\!\cdot\bm v_h)+S_h(\bm I_h\bm u,\bm v_h)\\
&=-a(\bm w,\bm v_h)-\gamma(\nabla\!\cdot\bm w,\nabla\!\cdot\bm v_h)+S_h(\bm I_h\bm u,\bm v_h).
\end{aligned}
\]
Substituting and rearranging,
\begin{equation}\label{eq:err}
\begin{aligned}
&(\CD\bm\zeta,\bm v_h)+a_h(\bm\zeta,\bm v_h)+b(\bm v_h,p_h-p)\\
&\qquad=(\CD\bm w,\bm v_h)+a(\bm w,\bm v_h)+\gamma(\nabla\!\cdot\bm w,\nabla\!\cdot\bm v_h)-S_h(\bm I_h\bm u,\bm v_h).
\end{aligned}
\end{equation}

By (F1) and $\nabla\!\cdot\bm u=0$ we have $b(\bm I_h\bm u,q_h)=b(\bm u,q_h)-b(\bm w,q_h)=0$ for all $q_h\in Q_h$. Together with $b(\bm u_h,q_h)=0$ this gives
\[
b(\bm\zeta,q_h)=b(\bm u_h,q_h)-b(\bm I_h\bm u,q_h)=0\qquad\text{for all }q_h\in Q_h .
\]
Since $p_h\in Q_h$, $b(\bm\zeta,p_h)=0$; and since $(\nabla\!\cdot\bm\zeta,q_h)=-b(\bm\zeta,q_h)=0$ for the chosen $q_h$,
\begin{equation}\label{eq:zeta-pressure}
b(\bm\zeta,p_h-p)=-b(\bm\zeta,p)=(\nabla\!\cdot\bm\zeta,p)=(\nabla\!\cdot\bm\zeta,p-q_h).
\end{equation}

Choose $\bm v_h=\bm\zeta$ in \eqref{eq:err}. The positivity result in \cref{lem:posc} gives $(\CD\bm\zeta,\bm\zeta)\ge\frac12\CD\|\bm\zeta\|^2$. This applies because $\bm\zeta\in W^{1,1}(0,T;L^2)$, as follows from the well-posedness discussion after \eqref{eq:semi} and $\bm I_h\bm u\in W^{1,1}(0,T;\bm V_h)$. By \cref{lem:bounds}(i), $a_h(\bm\zeta,\bm\zeta)=\nrm{\bm\zeta}^2$. Therefore,
\[
\tfrac12\CD\|\bm\zeta\|^2+\nrm{\bm\zeta}^2\le(\CD\bm w,\bm\zeta)+a(\bm w,\bm\zeta)+\gamma(\nabla\!\cdot\bm w,\nabla\!\cdot\bm\zeta)-S_h(\bm I_h\bm u,\bm\zeta)-(\nabla\!\cdot\bm\zeta,p-q_h).
\]
The terms on the right-hand side are bounded by using the Cauchy-Schwarz inequality and the definition of $\nrm{\cdot}$:
\[
\begin{aligned}
(\CD\bm w,\bm\zeta)&\le\|\CD\bm w\|\,\|\bm\zeta\|,\\
a(\bm w,\bm\zeta) &\le \nu^{1/2}\|\nabla\bm w\|\,\nrm{\bm\zeta} + \|\bmb\|_{L^\infty}\|\nabla\bm w\|\,\|\bm\zeta\|+\mu\|\bm w\|\,\|\bm\zeta\|,\\
\gamma(\nabla\!\cdot\bm w,\nabla\!\cdot\bm\zeta)+
S_h(\bm I_h\bm u,\bm\zeta)&\le \gamma^{1/2}\|\nabla\!\cdot\bm w\|\,\nrm{\bm\zeta} + S_h(\bm I_h\bm u,\bm I_h\bm u)^{1/2}\nrm{\bm\zeta},\\
-(\nabla\!\cdot\bm\zeta,p-q_h)&\le\|\nabla\!\cdot\bm\zeta\|\,\|p-q_h\|\le\gamma^{-1/2}\nrm{\bm\zeta}\,\|p-q_h\| .
\end{aligned}
\]
% The last estimate uses the grad-div part of the norm, $\gamma^{1/2}\|\nabla\!\cdot\bm\zeta\|\le\nrm{\bm\zeta}$. 
In the convection term, the factor $\|\bm\zeta\|$ is deliberately left to the Gronwall argument. Collecting the terms containing $\|\bm\zeta\|$ and those containing $\nrm{\bm\zeta}$, with $A$ and $B$ as in \cref{lem:interp},
\begin{equation}\label{eq:AB-bound}
\tfrac12\CD\|\bm\zeta\|^2+\nrm{\bm\zeta}^2\le A\|\bm\zeta\|+B\nrm{\bm\zeta} .
\end{equation}
By Young's inequality, 
$$A\|\bm\zeta\|\le\frac12A^2+\frac12\|\bm\zeta\|^2\; \text{and}\; B\nrm{\bm\zeta}\le\frac12B^2+\frac12\nrm{\bm\zeta}^2,$$ so
\[
\tfrac12\CD\|\bm\zeta\|^2+\tfrac12\nrm{\bm\zeta}^2\le\tfrac12\|\bm\zeta\|^2+\tfrac12\bigl(A^2+B^2\bigr).
\]
Multiplying by $2$ and using \cref{lem:interp}(ii), and (iii),
\begin{equation}\label{eq:G}
\CD\|\bm\zeta\|^2+\nrm{\bm\zeta}^2\le\|\bm\zeta\|^2+G,\qquad\text{with}\;
\quad G:=A^2+B^2\le C\bigl(\kappa_h^2+\vartheta_h^2\bigr).
\end{equation}
We will reuse \eqref{eq:G} in \cref{sec:full}. For the present estimate, return to \eqref{eq:AB-bound} and apply Young's inequality only to the $B$ term, $B\nrm{\bm\zeta}\le\frac12B^2+\frac12\nrm{\bm\zeta}^2$. This gives
\[
\CD\|\bm\zeta\|^2+\nrm{\bm\zeta}^2\le2A\|\bm\zeta\|+B^2 .
\]
This is the hypothesis of \cref{lem:sup} with $y=\|\bm\zeta\|$ (continuous, with $\|\bm\zeta\|^2$ absolutely continuous since $\bm\zeta\in W^{1,1}(0,T;L^2)$), $z=\nrm{\bm\zeta}^2$, $a=A$ and $g=B^2$. With $\bar A(t):=\sup_{s\le t}\RL A(s)$ and $\bar B(t):=\sup_{s\le t}\RL B^2(s)$, \eqref{eq:sup} yields
\begin{equation}\label{eq:zeta-final}
\sup_{0\le s\le t}\|\bm\zeta(s)\|\le\|\bm\zeta(0)\|+2\bar A(t)+\bar B(t)^{1/2},\qquad
\bigl(\RL\nrm{\bm\zeta}^2(t)\bigr)^{1/2}\le\|\bm\zeta(0)\|+\bar B(t)^{1/2}+\bigl(2\bar A(t)\sup_{s\le t}\|\bm\zeta(s)\|\bigr)^{1/2},
\end{equation}
and the last term is at most $\bar A(t)+\frac12\sup_{s\le t}\|\bm\zeta(s)\|$ by Young's inequality. By \cref{lem:interp}(ii),(iii), $\bar A\le C\sup_{s\le t}\RL(\kappa_h+\vartheta_h)(s)$ and $\bar B\le C\sup_{s\le t}\RL\vartheta_h^2(s)$.
The decomposition $\bm u_h-\bm u=\bm\zeta-\bm w$, the triangle inequality, and \cref{lem:interp}(i) (including $\RL\nrm{\bm w}^2\le C\RL\vartheta_h^2$) now give \eqref{eq:err-main}.
\end{proof}
\begin{proof}[Alternative proof via \cref{lem:gronwall}]
For comparison with \cref{sec:full} we record the Gronwall route, which yields the same rates with a Mittag-Leffler constant.
% Dropping $\nrm{\bm\zeta}^2\ge0$ in \eqref{eq:G} gives $\CD\|\bm\zeta\|^2\le\|\bm\zeta\|^2+G$. \eqref{lem:gronwall} with $y=\|\bm\zeta\|^2$, $c_1=1$ yields
Dropping the nonnegative term $\nrm{\bm\zeta}^2$ from \eqref{eq:G} and applying \cref{lem:gronwall} with $y=\|\bm\zeta\|^2$ and $c_1=1$ gives
\[
\|\bm\zeta(t)\|^2\le d\bigl(\|\bm\zeta(0)\|^2+\RL G\bigr),\qquad
\RL\|\bm\zeta\|^2\le t^\alpha d\Bigl(\frac{\|\bm\zeta(0)\|^2}{\Gamma(\alpha+1)}+\frac{\Gamma(\alpha)}{\Gamma(2\alpha)}\RL G\Bigr).
\]
Now apply $\RL$ to \eqref{eq:G} itself and use $\RL\CD y=y-y(0)$:
\[
\|\bm\zeta(t)\|^2-\|\bm\zeta(0)\|^2+\RL\nrm{\bm\zeta}^2\le\RL\|\bm\zeta\|^2+\RL G .
\]
Inserting the bound for $\RL\|\bm\zeta\|^2$,
\[
\|\bm\zeta(t)\|^2+\RL\nrm{\bm\zeta}^2\le\Bigl(1+\frac{T^\alpha d}{\Gamma(\alpha+1)}\Bigr)\|\bm\zeta(0)\|^2+\Bigl(1+\frac{T^\alpha d\,\Gamma(\alpha)}{\Gamma(2\alpha)}\Bigr)\RL G ,
\]
and therefore
\begin{equation}\label{eq:zeta-final-ML}
\|\bm\zeta(t)\|^2+\RL\nrm{\bm\zeta}^2\le K\bigl(\|\bm\zeta(0)\|^2+\RL G\bigr),\qquad
K:=1+T^\alpha d\max\Bigl\{\frac1{\Gamma(\alpha+1)},\frac{\Gamma(\alpha)}{\Gamma(2\alpha)}\Bigr\}.
\end{equation}
From $\bm u_h-\bm u=\bm\zeta-\bm w$ and $(a+b)^2\le2a^2+2b^2$ (for norms, also for seminorms),
\begin{equation}\label{eq:split}
\|\bm u_h-\bm u\|^2+\RL\nrm{\bm u_h-\bm u}^2\le2\Bigl(\|\bm\zeta\|^2+\RL\nrm{\bm\zeta}^2\Bigr)+2\Bigl(\|\bm w\|^2+\RL\nrm{\bm w}^2\Bigr).
\end{equation}
By \eqref{eq:zeta-final-ML}, the first bracket in \eqref{eq:split} is bounded by $K(\|\bm\zeta(0)\|^2+\RL G)$, where $\|\bm\zeta(0)\|=\|\bm u_{h,0}-\bm I_h\bm u_0\|$ and $G\le C(\kappa_h^2+\vartheta_h^2)$ by \eqref{eq:G}. The interpolation estimate in \cref{lem:interp}(i) bounds the second bracket by $C(\vartheta_h(t)^2+\RL\vartheta_h^2)$. Consequently,
\[
\|\bm u_h(t)-\bm u(t)\|^2+\RL\nrm{\bm u_h-\bm u}^2\le C_{T,\alpha}\Bigl[\|\bm u_{h,0}-\bm I_h\bm u_0\|^2+\vartheta_h(t)^2+\RL\bigl(\vartheta_h^2+\kappa_h^2\bigr)\Bigr],
\]
Here $C_{T,\alpha}=C(C_I,C_S)K$; the estimate is pointwise in $t$ and includes the Mittag--Leffler factor $d$.
\end{proof}

\begin{corollary}\label{cor:rate}
For $\bm u_{h,0}=\bm I_h\bm u_0$,
\[
\begin{aligned}
&\sup_t\|\bm u_h-\bm u\|+\bigl(\RL\nrm{\bm u_h-\bm u}^2\bigr)^{1/2}\\
&\qquad\le C_{T,\alpha}\Bigl[c_hh^r\sup_t\|\bm u\|_{r+1}+\gamma^{-1/2}h^{s+1}\sup_t\|p\|_{s+1}+h^{r+1}\sup_t\|\CD\bm u\|_{r+1}\Bigr],
\end{aligned}
\]
The constant is $C_{T,\alpha}=C(C_I,C_S)\bigl(1+T^\alpha/\Gamma(\alpha+1)\bigr)$. For Taylor--Hood ($s=r-1$) and $\gamma=O(1)$, the estimate is $O(h^r)$.
\end{corollary}
\begin{proof}
By \eqref{eq:rl},
\[
\RL(\kappa_h+\vartheta_h)(s)\le\frac{T^\alpha}{\Gamma(\alpha+1)}\sup_t(\kappa_h+\vartheta_h),\qquad
\RL\vartheta_h^2(s)\le\frac{T^\alpha}{\Gamma(\alpha+1)}\sup_t\vartheta_h^2 ,
\]
and $(T^\alpha/\Gamma(\alpha+1))^{1/2}\le1+T^\alpha/\Gamma(\alpha+1)$. Substituting these bounds into \eqref{eq:err-main} proves the result.
\end{proof}

\begin{remark}[On the hypotheses]\label{rem:hyp}
\begin{enumerate}[(i),leftmargin=2.4em]
\item \emph{How the assumptions enter.} The Fortin property (F1) gives $b(\bm\zeta,q_h)=0$, while (F2) and (S2) are used in \cref{lem:interp}. The bound (S1) is not needed for \cref{thm:err}; it enters through \cref{lem:bounds}(ii) in the pressure estimates. The analysis uses only $\|\bmb\|_{L^\infty}$ and $\nabla\!\cdot\bmb=0$ from the advection field; the stronger assumption $\bmb\in W^{1,\infty}$ is needed only to bound the constants in (S1)--(S2). If $\nabla\!\cdot\bmb\ne0$, replace $\mu\|\bm v\|^2$ by $\bigl((\mu-\tfrac12\nabla\!\cdot\bmb)\bm v,\bm v\bigr)$ and assume $\mu-\tfrac12\nabla\!\cdot\bmb\ge0$.
\item \emph{The $L^2$ approximation property (F3) is optional for the rate $h^r$.} It is used only to bound $\|\CD\bm w\|$ and $\mu\|\bm w\|$. Without (F3), use $\|\CD\bm w\|\le C_FC_Ih^r\|\CD\bm u\|_{r+1}$ and $\mu\|\bm w\|\le\mu C_FC_Ih^r\|\bm u\|_{r+1}$ instead. Then \cref{cor:rate} remains valid with $\kappa_h$ replaced by $h^r\|\CD\bm u\|_{r+1}$ and with $\mu h+\mu^{1/2}h$ in $c_h$ replaced by $\mu C_F$. Conversely, if (F3) holds in its natural form $\|\bm v-\bm I_h\bm v\|\le C_Ih^m\|\bm v\|_m$ for $1\le m\le r+1$, it is enough to assume $\CD\bm u\in L^\infty(0,T;H^{r})$ in \eqref{sec:error}.
\item \emph{Sharpness.} For Taylor--Hood elements, the pressure approximation term $\gamma^{-1/2}h^{s+1}=\gamma^{-1/2}h^r$ limits the rate to $h^r$, regardless of how convection is treated. This agrees with the integer-order result in \cite{dFGJN}. Stabilised Oseen methods can achieve the higher order $h^{r+1/2}$ \cite{BF07}, but doing so here would require an interpolant orthogonal to the fluctuation space of $S_h$. That improvement is useful only when $\nabla\!\cdot\bm V_h\subset Q_h$, so the pressure term vanishes.
\item \emph{Dimension.} The proof extends directly to three space dimensions; replace $\sqrt2$ by $\sqrt3$ in \cref{lem:bounds}.
\end{enumerate}
\end{remark}

\begin{proposition}[Pressure error]\label{prop:pressure}
For almost every $t$,
\begin{equation}\label{eq:p-cond}
\beta_1\|p_h-q_h\|\le C_F\|\CD\bm\zeta\|+M_h\nrm{\bm\zeta}+\|\bmb\|_{L^\infty}\|\bm\zeta\|+\Theta_h,
\end{equation}
where
\[
\begin{aligned}
\Theta_h:={}&C_F\|\CD\bm w\|+(\nu+2\gamma+C_F\|\bmb\|_{L^\infty})\|\nabla\bm w\|+C_F\mu\|\bm w\|+(C_Sh)^{1/2}S_h(\bm I_h\bm u,\bm I_h\bm u)^{1/2}\\
&+\sqrt2\|p-q_h\| .
\end{aligned}
\]
The fractionally integrated pressure also admits an unconditional estimate:
\begin{equation}\label{eq:p-int}
\begin{aligned}
\beta_1\|\RL(p_h-q_h)(t)\|\le{}&C_F\bigl(\|\bm\zeta(t)\|+\|\bm\zeta(0)\|+\|\bm w(t)\|+\|\bm w(0)\|\bigr)\\
&+M_h\RL\nrm{\bm\zeta}+\|\bmb\|_{L^\infty}\RL\|\bm\zeta\|+\RL\Theta_h',
\end{aligned}
\end{equation}
Set $\Theta_h'=\Theta_h-C_F\|\CD\bm w\|$. The right-hand side of \eqref{eq:p-int} is bounded using \cref{thm:err} and \cref{lem:interp}. In particular, for $\bm u_{h,0}=\bm I_h\bm u_0$ and Taylor--Hood elements, $\beta_1\|\RL(p_h-q_h)(t)\|\le C(1+T^\alpha)h^r$, with $C$ independent of $\nu^{-1}$.
\end{proposition}
\begin{proof}
Since $p_h-q_h\in Q_h$, the inf--sup condition \eqref{eq:infsup} gives
\[
\beta_1\|p_h-q_h\|\le\sup_{0\ne\bm v\in\bm V_h}\frac{b(\bm v,p_h-q_h)}{\|\nabla\bm v\|}.
\]
By \eqref{eq:err}, for every $\bm v\in\bm V_h$,
\[
\begin{aligned}
b(\bm v,p_h-q_h)&=b(\bm v,p_h-p)+b(\bm v,p-q_h)\\
&=-(\CD\bm\zeta,\bm v)-a_h(\bm\zeta,\bm v)+(\CD\bm w,\bm v)+a(\bm w,\bm v)\\
&\quad+\gamma(\nabla\!\cdot\bm w,\nabla\!\cdot\bm v)-S_h(\bm I_h\bm u,\bm v)+b(\bm v,p-q_h).
\end{aligned}
\]
We bound the terms by $\|\nabla\bm v\|$, using $\|\bm v\|\le C_F\|\nabla\bm v\|$ and \cref{lem:bounds}:
\[
\begin{aligned}
|(\CD\bm\zeta,\bm v)|&\le C_F\|\CD\bm\zeta\|\,\|\nabla\bm v\|,\\
|a_h(\bm\zeta,\bm v)|&\le\bigl(M_h\nrm{\bm\zeta}+\|\bmb\|_{L^\infty}\|\bm\zeta\|\bigr)\|\nabla\bm v\| &&\text{by \eqref{eq:ah-bound}},\\
|(\CD\bm w,\bm v)|&\le C_F\|\CD\bm w\|\,\|\nabla\bm v\|,\\
|a(\bm w,\bm v)+\gamma(\nabla\!\cdot\bm w,\nabla\!\cdot\bm v)|&\le\bigl[(\nu+2\gamma+C_F\|\bmb\|_{L^\infty})\|\nabla\bm w\|\\
&\qquad+C_F\mu\|\bm w\|\bigr]\|\nabla\bm v\| &&\text{by \eqref{eq:a-bound}},\\
|S_h(\bm I_h\bm u,\bm v)|&\le S_h(\bm I_h\bm u,\bm I_h\bm u)^{1/2}(C_Sh)^{1/2}\|\nabla\bm v\|,\\
|b(\bm v,p-q_h)|&\le\sqrt2\,\|\nabla\bm v\|\,\|p-q_h\| &&\text{by \eqref{eq:b-bound}}.
\end{aligned}
\]
Adding these bounds and dividing by $\|\nabla\bm v\|$ gives \eqref{eq:p-cond}.

For \eqref{eq:p-int}, fix $\bm v\in\bm V_h$ and apply $\RL$ to the identity above. Since $\bm\zeta$, $\bm w$ are absolutely continuous, $\RL\CD\bm\zeta=\bm\zeta-\bm\zeta(0)$ and $\RL\CD\bm w=\bm w-\bm w(0)$, hence
\[
\begin{aligned}
b\bigl(\bm v,\RL(p_h-q_h)\bigr)={}&-(\bm\zeta(t)-\bm\zeta(0),\bm v)-a_h(\RL\bm\zeta,\bm v)+(\bm w(t)-\bm w(0),\bm v)\\
&+\RL\Bigl[a(\bm w,\bm v)+\gamma(\nabla\!\cdot\bm w,\nabla\!\cdot\bm v)-S_h(\bm I_h\bm u,\bm v)+b(\bm v,p-q_h)\Bigr].
\end{aligned}
\]
By \eqref{eq:ah-bound} and $\nrm{\RL\bm\zeta}\le\RL\nrm{\bm\zeta}$, $\|\RL\bm\zeta\|\le\RL\|\bm\zeta\|$,
\[
|a_h(\RL\bm\zeta,\bm v)|\le\bigl(M_h\RL\nrm{\bm\zeta}+\|\bmb\|_{L^\infty}\RL\|\bm\zeta\|\bigr)\|\nabla\bm v\| .
\]
The remaining terms are bounded as before. Since $\RL(p_h-q_h)\in Q_h$, dividing by $\|\nabla\bm v\|$ and taking the supremum gives \eqref{eq:p-int}. Finally, $\sup_{s\le t}\|\bm\zeta(s)\|$ and $\RL\nrm{\bm\zeta}^2$ are bounded in \eqref{eq:zeta-final}, $\RL\|\bm\zeta\|\le(t^\alpha/\Gamma(\alpha+1))\sup_{s\le t}\|\bm\zeta(s)\|$, and $\RL\nrm{\bm\zeta}\le(t^\alpha/\Gamma(\alpha+1))^{1/2}(\RL\nrm{\bm\zeta}^2)^{1/2}$ by \eqref{eq:rl}.
\end{proof}
The pointwise estimate \eqref{eq:p-cond} is \emph{conditional}: it involves $\|\CD\bm\zeta\|$, which the energy argument does not control without introducing a negative power of $\nu$; see the discussion after \cref{thm:semidiscrete_pressure_stability}. By contrast, the integrated estimate \eqref{eq:p-int} is viscosity robust. As in the integer-order case \cite{BF07}, it is the time-integrated pressure that can be bounded uniformly in $\nu$.

\begin{remark}\label{rem:variants}
\begin{enumerate}[(a),leftmargin=2.4em]
\item \emph{Reaction term.} When $\mu>0$, the reaction term controls the $L^2$ norm and lets us avoid the Gronwall lemma. Indeed, $A\|\bm\zeta\|\le A\mu^{-1/2}\nrm{\bm\zeta}$, so that
\[
\tfrac12\CD\|\bm\zeta\|^2+\nrm{\bm\zeta}^2\le(\mu^{-1/2}A+B)\nrm{\bm\zeta}\le\tfrac12\nrm{\bm\zeta}^2+\tfrac12(\mu^{-1/2}A+B)^2,
\]
and hence
\[
\CD\|\bm\zeta\|^2+\nrm{\bm\zeta}^2\le2\mu^{-1}A^2+2B^2,
\]
and, applying $\RL$,
\[
\|\bm\zeta(t)\|^2+\RL\nrm{\bm\zeta}^2\le\|\bm\zeta(0)\|^2+2\,\RL\bigl(\mu^{-1}A^2+B^2\bigr).
\]
The constants contain $\mu^{-1}$ but never $\nu^{-1}$.
\item \emph{Grad-div parameter.} The assumption $\gamma>0$ is needed to bound the pressure approximation term $\gamma^{-1/2}\|p-q_h\|$. If $\gamma=0$, we only have
\[
-(\nabla\!\cdot\bm\zeta,p-q_h)\le\sqrt2\|\nabla\bm\zeta\|\,\|p-q_h\|\le\sqrt2\,\nu^{-1/2}\nrm{\bm\zeta}\,\|p-q_h\|,
\]
so that $\nu^{-1/2}$ multiplies the pressure approximation error, unless $\nabla\!\cdot\bm V_h\subset Q_h$ (then $\nabla\!\cdot\bm\zeta=0$ and the term vanishes). For Taylor--Hood, $\gamma\sim1$ balances $\gamma^{1/2}h^r$ and $\gamma^{-1/2}h^{s+1}$.
\item \emph{Convection.} In our estimate, convection contributes $\|\bmb\|_{L^\infty}h^r$. Moving the derivative to the error instead gives
\[
\nu^{-1/2}\|\bmb\|_{L^\infty}\|\bm w\|\,\nrm{\bm\zeta}\lesssim\nu^{-1/2}\|\bmb\|_{L^\infty}h^{r+1}\,\nrm{\bm\zeta},
\]
which reintroduces $\nu^{-1/2}$ (at higher order); we do not use it.
\end{enumerate}
\end{remark}
% =====================================================================
\section{Fully discrete schemes}
\label{sec:full}

%\subsection{The uniform L1 scheme}

In this section we analyse a fully discrete scheme based on L1 method \cite{LX}.
Divide the time interval $[0,T]$ into $N$ equal subintervals with a time step size
$\tau=T/N$. Let $t_n=n\tau$, $\varphi^n=\varphi(t_n)$, and
$\kappa=\tau^\alpha$.
The L1 approximation of the Caputo fractional derivative of $\varphi(t)$ at time $t_n$ is defined by
\begin{equation}
\label{eq:Dt}
\Dt \varphi^n
=
\tau^{-\alpha}\sum_{k=1}^n
\beta_{n-k}(\varphi^k-\varphi^{k-1}),
\end{equation}
where
$
\beta_j
=
[(j+1)^{1-\alpha}-j^{1-\alpha}]/\Gamma(2-\alpha),
\, j\ge0.
$

The fully discrete scheme reads as follows:
given $\bm U_h^0=\bm u_{h,0}$, find
$(\bm U_h^n,P_h^n)\in\bm V_h\times Q_h$,
$1\le n\le N$, such that
\begin{equation}
\label{eq:fd}
\begin{aligned}
(\Dt\bm U_h^n,\bm v_h)
+a_h(\bm U_h^n,\bm v_h)
+b(\bm v_h,P_h^n)
&=(\bm f^n,\bm v_h),\\
b(\bm U_h^n,q_h)&=0
\end{aligned}
\end{equation}
for all $(\bm v_h,q_h)\in\bm V_h\times Q_h$.
At each time level, the resulting saddle-point problem
is well posed by the coercivity of the bilinear form
\[
\beta_0\tau^{-\alpha}(\bm w_h,\bm v_h)
+a_h(\bm w_h,\bm v_h)
\]
on the discrete divergence-free subspace and the
discrete inf-sup condition.


Let $\{\chi_j\}_{j\ge0}$ denote the complementary
L1 weights defined by
\[
\sum_{j=k}^n\chi_{n-j}\beta_{j-k}=1,
\qquad 1\le k\le n.
\]
We recall the following standard properties of the L1
and complementary weights; see, e.g., [--].

\begin{lemma}[Properties of the L1 weights]
\label{lem:weights}
The L1 weights $\beta_j$ are positive and decreasing.
Moreover, for any sequence $\{\bm W^n\}$ in a real
Hilbert space,
\[
2(\Dt\bm W^n,\bm W^n)
\ge
\Dt\|\bm W^n\|^2.
\]
The complementary weights satisfy $\chi_j\ge0$ and
\[
\kappa\sum_{j=1}^n\chi_{n-j}\Dt y^j
=
y^n-y^0
\]
for any scalar sequence $\{y^n\}$. In addition,
\[
\kappa\sum_{j=0}^{n-1}\chi_j
\le Ct_n^\alpha,
\qquad 1\le n\le N.
\]
\end{lemma}


%------------------------



\begin{lemma}[Discrete fractional Gronwall inequality]
\label{lem:dgronwall}
Let $Z^n,G^n\ge0$ and $A\ge0$ satisfy
\[
Z^n\le A+\kappa\sum_{j=1}^n
\chi_{n-j}(Z^j+G^j),
\qquad 1\le n\le N.
\]
If $\kappa\chi_0\le 1/2$, then
\[
Z^n\le C_{T,\alpha}\left(
A+\max_{1\le m\le n}
\kappa\sum_{j=1}^m\chi_{m-j}G^j
\right),
\qquad 1\le n\le N,
\]
where $C_{T,\alpha}$ is independent of $h$ and $\tau$.
\end{lemma}

%--------------------

\begin{comment}

\begin{proof}
Fix $n\le N$ and set
\[
D=A+\max_{1\le m\le n}
\kappa\sum_{j=1}^m\chi_{m-j}G^j.
\]
Then
\begin{equation}\label{eq:dg-basic}
Z^m\le D+\kappa\sum_{j=1}^m
\chi_{m-j}Z^j,\qquad 1\le m\le n.
\end{equation}
By Lemma~\ref{lem:weights},
\[
\kappa\sum_{j=0}^{\ell-1}\chi_j
\le C_\alpha(\ell\tau)^\alpha.
\]
Choose $\delta>0$ such that
$C_\alpha\delta^\alpha\le1/2$, and partition
$\{1,\ldots,n\}$ into consecutive blocks of length
at most
\[
L=\max\{1,\lfloor\delta/\tau\rfloor\}.
\]
For any block $\{a,\ldots,b\}$ and $a\le m\le b$,
\begin{equation}\label{eq:dg-block}
\kappa\sum_{j=a}^m\chi_{m-j}\le\frac12.
\end{equation}
Indeed, this follows from the cumulative weight
bound when $\tau\le\delta$, and from
$\kappa\chi_0\le1/2$ when $\tau>\delta$.

Let $M_k$ denote the maximum of $Z^j$ over the
first $k$ blocks, with $M_0=0$.
For the first block, \eqref{eq:dg-basic} and
\eqref{eq:dg-block} give
\[
M_1\le 2D.
\]
For a subsequent block $\{a,\ldots,b\}$,
splitting the convolution at $a$ yields
\[
Z^m
\le D+C_\alpha T^\alpha M_{k-1}
+\frac12 M_k,
\qquad a\le m\le b.
\]
Taking the maximum over the current block
and combining it with the bound for the
preceding blocks, we obtain
\[
M_k\le
2D+(1+2C_\alpha T^\alpha)M_{k-1}.
\]
Consequently,
\[
M_k\le
2D\sum_{\ell=0}^{k-1}
(1+2C_\alpha T^\alpha)^\ell.
\]
The number of blocks is bounded by
$1+2T/\delta$, independently of $\tau$.
Hence $M_k\le C_{T,\alpha}D$, which proves
the assertion.
\end{proof}


\end{comment}

%-------------------------



\begin{theorem}[Fully discrete stability]
\label{thm:fdstab}
The solution of \eqref{eq:fd} satisfies
\[
\frac12\Dt\|\bm U_h^n\|^2
+\nrm{\bm U_h^n}^2
\le(\bm f^n,\bm U_h^n),
\qquad 1\le n\le N.
\]
If $\kappa\chi_0\le1/2$, then
\[
\|\bm U_h^n\|^2
+\kappa\sum_{j=1}^n
\chi_{n-j}\nrm{\bm U_h^j}^2
\le C_{T,\alpha}
\left(
\|\bm U_h^0\|^2
+\max_{1\le m\le n}
\kappa\sum_{j=1}^m
\chi_{m-j}\|\bm f^j\|^2
\right).
\]
%where $C_{T,\alpha}$ is independent of $h$ and $\tau$.
If $\mu>0$, then, for every $\tau>0$,
\[
\|\bm U_h^n\|^2
+\kappa\sum_{j=1}^n
\chi_{n-j}\nrm{\bm U_h^j}^2\le
\|\bm U_h^0\|^2
+\mu^{-1}\kappa\sum_{j=1}^n
\chi_{n-j}\|\bm f^j\|^2.
\]
\end{theorem}

\begin{proof}
Taking $\bm v_h=\bm U_h^n$ and $q_h=P_h^n$
in \eqref{eq:fd}, the discrete incompressibility
constraint and Lemma~\ref{lem:weights} give the first estimate.
By Young's inequality,
\begin{equation}\label{eq:fdstab-energy}
\Dt Z^n+2\nrm{\bm U_h^n}^2
\le Z^n+G^n,
\end{equation}
where
$Z^n=\|\bm U_h^n\|^2$ and $G^n=\|\bm f^n\|^2.$
Multiplying \eqref{eq:fdstab-energy} at time level
$j$ by $\kappa\chi_{n-j}$, summing over
$j=1,\ldots,n$, and using Lemma~\ref{lem:weights},
we obtain
\begin{equation}\label{eq:fdstab-conv}
Z^n+2\kappa\sum_{j=1}^n
\chi_{n-j}\nrm{\bm U_h^j}^2
\le Z^0+\kappa\sum_{j=1}^n
\chi_{n-j}(Z^j+G^j).
\end{equation}
Dropping the nonnegative energy term and applying
Lemma~\ref{lem:dgronwall} yields
\begin{equation}\label{eq:fdstab-Z}
Z^n\le C_{T,\alpha}
\left(
Z^0+\max_{1\le m\le n}
\kappa\sum_{j=1}^m\chi_{m-j}G^j
\right).
\end{equation}
Since this estimate holds at every time level
$j\le n$, the positivity of $\chi_j$ and
Lemma~\ref{lem:weights} imply
\[
\begin{aligned}
\kappa\sum_{j=1}^n\chi_{n-j}Z^j
&\le C_{T,\alpha}
\left(
Z^0+\max_{1\le m\le n}
\kappa\sum_{k=1}^m\chi_{m-k}G^k
\right)
\kappa\sum_{j=0}^{n-1}\chi_j\\
&\le C_{T,\alpha}
\left(
Z^0+\max_{1\le m\le n}
\kappa\sum_{k=1}^m\chi_{m-k}G^k
\right).
\end{aligned}
\]
Substituting this bound and \eqref{eq:fdstab-Z}
into \eqref{eq:fdstab-conv} proves the second
assertion.
%
If $\mu>0$, then
$\nrm{\bm U_h^n}^2\ge\mu\|\bm U_h^n\|^2$.
Consequently,
\[
(\bm f^n,\bm U_h^n)
\le\mu^{-1/2}\|\bm f^n\|
\nrm{\bm U_h^n}\\
\le\frac12\nrm{\bm U_h^n}^2
+\frac1{2\mu}\|\bm f^n\|^2.
\]
The first stability inequality therefore gives
$
\Dt Z^n+\nrm{\bm U_h^n}^2
\le\mu^{-1}G^n.
$
Multiplying by $\kappa\chi_{n-j}$ at time
level $j$, summing over $j=1,\ldots,n$, and
using the complementary-weight identity in
Lemma~\ref{lem:weights} yields
\[
Z^n+\kappa\sum_{j=1}^n
\chi_{n-j}\nrm{\bm U_h^j}^2
\le Z^0+\mu^{-1}\kappa\sum_{j=1}^n
\chi_{n-j}G^j.
\]
This completes the proof.
\end{proof}

%-------------------

\begin{comment}

\begin{remark}[Unconditional stability]
\label{rem:unconditional-stability}
Under homogeneous Dirichlet boundary conditions, the
Poincar\'e inequality implies
\[
\nrm{\bm v_h}^2
\ge \lambda_0\|\bm v_h\|^2,
\qquad
\lambda_0=\mu+\frac{\nu}{C_P^2}>0.
\]
Consequently, the argument used for $\mu>0$ in
Theorem~\ref{thm:fdstab} yields
\[
\begin{aligned}
&\|\bm U_h^n\|^2
+\kappa\sum_{j=1}^n
\chi_{n-j}\nrm{\bm U_h^j}^2\\
&\qquad\le
\|\bm U_h^0\|^2
+\lambda_0^{-1}\kappa\sum_{j=1}^n
\chi_{n-j}\|\bm f^j\|^2.
\end{aligned}
\]
Thus, the fully discrete scheme is unconditionally
stable for every $\mu\ge0$, without any restriction
on the time step.
\end{remark}


\end{comment}
%--------------------

\paragraph{Temporal consistency.}
Let $\mathcal R^n$ denote the consistency error
of the L1 approximation,
\[
\mathcal R^n
:=\Dt\bm u^n-\CD\bm u(t_n),
\qquad
\bm u^n=\bm u(t_n).
\]
We recall the standard error estimate of the L1 method; see \cite{LX}.

\begin{lemma}
\label{lem:cons}
Let $X$ be a Banach space and
$u\in C^2([0,T];X)$. Then
\[
\|\Dt u^n-\CD u(t_n)\|_X
\le
C\tau^{2-\alpha}
\max_{t\in[0,T]}
\|\partial_t^2 u(t)\|_X,
\qquad 1\le n\le N,
\]
where $C>0$ may depend on $\alpha$, but is
independent of $\tau$ and $n$.
\end{lemma}


%----------------------------------

\begin{comment}

\begin{lemma}[Increment of the L1 consistency error]
\label{lem:L1_increment}
Let $X$ be a Hilbert space and let $0<\alpha<1$.
Suppose that
\[
u\in C^3([0,T];X),\qquad
u^{(4)}\in L^\infty(0,T;X),
\qquad u''(0)=0.
\]
Let
\[
\mathcal R^n
=D_\tau^\alpha u^n-{}^CD_t^\alpha u(t_n),
\qquad \mathcal R^0=0.
\]
Then, for $0\le n\le N-1$,
\[
\|\mathcal R^{n+1}-\mathcal R^n\|_X
\le C_{\alpha,T}\tau^{3-\alpha}
\left(
\|u^{(3)}\|_{L^\infty(0,T;X)}
+
\|u^{(4)}\|_{L^\infty(0,T;X)}
\right).
\]
\end{lemma}

\end{comment}

%-----------------------

We are now ready to derive the fully discrete velocity error estimate.

\begin{theorem}[Fully discrete velocity error]
\label{thm:fderr}
Assume the hypotheses of Section~\ref{sec:error}
and suppose that
$\tau^\alpha\Gamma(2-\alpha)\le 1/2$.
In addition, let
\[
\bm u\in
C^2\bigl([0,T];[H^{r+1}\cap H_0^1]^2\bigr).
\]
Then, for $1\le n\le N$,
\[
\begin{aligned}
&\|\bm U_h^n-\bm u^n\|^2
+\kappa\sum_{j=1}^n
\chi_{n-j}\nrm{\bm U_h^j-\bm u^j}^2\\
&\qquad\le C_{T,\alpha}\left[
\|\bm U_h^0-\bm I_h\bm u^0\|^2
+h^{2(r+1)}
+\tau^{2(2-\alpha)}
+\max_{1\le j\le n}
\vartheta_h(t_j)^2
\right],
\end{aligned}
\]
where $C_{T,\alpha}$ is independent of $h$,
$\tau$, and negative powers of $\nu$,
but may depend on the solution norms.
In particular, if
$\bm U_h^0=\bm I_h\bm u^0$, then
\[
\|\bm U_h^n-\bm u^n\|
=
O\bigl(
c_hh^r+\gamma^{-1/2}h^{s+1}
+h^{r+1}+\tau^{2-\alpha}
\bigr),
\qquad 1\le n\le N.
\]
%provided the corresponding solution norms remain bounded.
\end{theorem}



%--------------------------------




\begin{proof}
Let
$
\bm\zeta^n=\bm U_h^n-\bm I_h\bm u^n,$  
$\bm w^n=\bm u^n-\bm I_h\bm u^n,$
and $q_h^n\in Q_h$ be a best approximation
of $p^n$.
Since
$\CD\bm u(t_n)=\Dt\bm u^n-\mathcal R^n$,
subtracting the continuous problem at $t_n$
from \eqref{eq:fd} gives
\begin{equation}\label{eq:fd-err}
\begin{aligned}
&(\Dt\bm\zeta^n,\bm v_h)
+a_h(\bm\zeta^n,\bm v_h)
+b(\bm v_h,P_h^n-p^n)\\
&\qquad=
(\Dt\bm w^n-\mathcal R^n,\bm v_h)
+a(\bm w^n,\bm v_h)+
\gamma(\nabla\cdot\bm w^n,
       \nabla\cdot\bm v_h)
-S_h(\bm I_h\bm u^n,\bm v_h).
\end{aligned}
\end{equation}
%
By the discrete incompressibility constraint
and the Fortin property,
\[
b(\bm\zeta^n,q_h)=0
\qquad\forall q_h\in Q_h.
\]
Taking $\bm v_h=\bm\zeta^n$ in
\eqref{eq:fd-err}, using
Lemma~\ref{lem:interp}, Young's inequality,
and the discrete energy inequality in
Lemma~\ref{lem:weights}, we obtain
\begin{equation}\label{eq:fd-G}
\Dt\|\bm\zeta^n\|^2
+\nrm{\bm\zeta^n}^2
\le
\|\bm\zeta^n\|^2+G^n,
\end{equation}
where
$
G^n=A_n^2+B_n^2,
$
with
\[
\begin{aligned}
A_n={}&
\|\Dt\bm w^n\|
+\|\mathcal R^n\|
+\|\bmb\|_{L^\infty}\|\nabla\bm w^n\|
+\mu\|\bm w^n\|,
\\
B_n={}&C\bigl[
(\sqrt\nu+\sqrt\gamma+\sqrt h)
h^r\|\bm u^n\|_{r+1}
+\gamma^{-1/2}h^{s+1}\|p^n\|_{s+1}
\bigr].
\end{aligned}
\]
%
Since $\bm I_h$ is linear and independent of time,
$
\Dt\bm w^n=(I-\bm I_h)\Dt\bm u^n.
$
The approximation properties of $\bm I_h$
and the definition of the L1 operator imply
\[
\|\Dt\bm w^n\|
\le Ch^{r+1}\|\Dt\bm u^n\|_{r+1}
\le C_{T,\alpha}h^{r+1}
\max_{0\le t\le T}
\|\partial_t\bm u(t)\|_{r+1}.
\]
Moreover, Lemma~\ref{lem:cons} gives
\[
\|\mathcal R^n\|
\le C_\alpha\tau^{2-\alpha}
\max_{0\le t\le T}
\|\partial_t^2\bm u(t)\|.
\]
Consequently, by the regularity assumptions
and Lemma~\ref{lem:interp},
\begin{equation}\label{eq:fd-G-bound}
G^n\le C_{T,\alpha}
\left(
h^{2(r+1)}
+\tau^{2(2-\alpha)}
+\vartheta_h(t_n)^2
\right).
\end{equation}
%
Multiplying \eqref{eq:fd-G} at time level $j$
by $\kappa\chi_{n-j}$, summing over
$j=1,\ldots,n$, and using the complementary-weight
identity yields
\begin{equation}\label{eq:fd-weighted}
\|\bm\zeta^n\|^2
+\kappa\sum_{j=1}^n
\chi_{n-j}\nrm{\bm\zeta^j}^2\le
\|\bm\zeta^0\|^2
+\kappa\sum_{j=1}^n
\chi_{n-j}
\left(\|\bm\zeta^j\|^2+G^j\right).
\end{equation}
%
Dropping the nonnegative energy term
and applying Lemma~\ref{lem:dgronwall},
under the condition
$\tau^\alpha\Gamma(2-\alpha)\le1/2$,
we obtain
\begin{equation}\label{eq:fd-zeta-L2}
\|\bm\zeta^n\|^2
\le C_{T,\alpha}
\left(
\|\bm\zeta^0\|^2
+\max_{1\le m\le n}
\kappa\sum_{j=1}^m
\chi_{m-j}G^j
\right).
\end{equation}
%
By \eqref{eq:fd-G-bound}, the positivity
of the complementary weights, and
Lemma~\ref{lem:weights},
\[
\max_{1\le m\le n}
\kappa\sum_{j=1}^m\chi_{m-j}G^j
\le C_{T,\alpha}
\left(
h^{2(r+1)}
+\tau^{2(2-\alpha)}
+\max_{1\le j\le n}\vartheta_h(t_j)^2
\right).
\]
Hence,
\begin{equation}\label{eq:fd-zeta-bound}
\|\bm\zeta^n\|^2
\le C_{T,\alpha}\bigg[
\|\bm\zeta^0\|^2
+h^{2(r+1)}
+\tau^{2(2-\alpha)}
+\max_{1\le j\le n}
\vartheta_h(t_j)^2
\bigg].
\end{equation}
%
Applying \eqref{eq:fd-zeta-bound}
at each time level $j\le n$,
and using the positivity and bounded sum
of the complementary weights in
\eqref{eq:fd-weighted}, we conclude that
\begin{equation}\label{eq:fd-zeta-energy}
\begin{aligned}
&\|\bm\zeta^n\|^2
+\kappa\sum_{j=1}^n
\chi_{n-j}\nrm{\bm\zeta^j}^2\\
&\qquad\le C_{T,\alpha}\bigg[
\|\bm\zeta^0\|^2
+h^{2(r+1)}
+\tau^{2(2-\alpha)}
+\max_{1\le j\le n}
\vartheta_h(t_j)^2
\bigg].
\end{aligned}
\end{equation}
%
Finally, since
$
\bm U_h^n-\bm u^n
=\bm\zeta^n-\bm w^n,
$
the triangle inequality and
Lemma~\ref{lem:interp} give
\[
\begin{aligned}
\|\bm U_h^n-\bm u^n\|^2
&\le
2\|\bm\zeta^n\|^2+2\|\bm w^n\|^2,\\
\nrm{\bm U_h^j-\bm u^j}^2
&\le
2\nrm{\bm\zeta^j}^2
+2\nrm{\bm w^j}^2,
\end{aligned}
\]
where
$
\|\bm w^j\|^2+\nrm{\bm w^j}^2
\le C\vartheta_h(t_j)^2.
$
Combining these bounds with
\eqref{eq:fd-zeta-energy} and
Lemma~\ref{lem:weights}
proves the asserted estimate.

If $\bm U_h^0=\bm I_h\bm u^0$,
then $\bm\zeta^0=0$, and the stated
convergence rate follows from the
spatial approximation estimates.
\end{proof}
%----------------------------


\begin{proposition}[Fully discrete integrated pressure error]
\label{prop:fdpressure}
Assume the hypotheses of Theorem~\ref{thm:fderr}.
Then, for $1\le n\le N$,
\[
\left\|
\kappa\sum_{j=1}^n
\chi_{n-j}(P_h^j-p^j)
\right\|\le C_{T,\alpha,\beta_1}
(1+M_h)
\left(
\|\bm\zeta^0\|
+h^{r+1}
+\tau^{2-\alpha}
+\max_{1\le j\le n}\vartheta_h(t_j)
\right),
\]
where
$
\bm\zeta^0=\bm U_h^0-\bm I_h\bm u^0,
$
and $C_{T,\alpha,\beta_1}$ is independent
of $h$, $\tau$, and negative powers of $\nu$.
\end{proposition}


\begin{proof}
Let
\[
\widetilde\rho^n
=
\kappa\sum_{j=1}^n
\chi_{n-j}(P_h^j-q_h^j),
\qquad
\widetilde{\bm\zeta}^n
=
\kappa\sum_{j=1}^n
\chi_{n-j}\bm\zeta^j,
\]
where $q_h^j\in Q_h$ is a best approximation
of $p^j$.
Adding $b(\bm v_h,p^j-q_h^j)$ to both sides
of \eqref{eq:fd-err}, multiplying the resulting
identity by $\kappa\chi_{n-j}$, and summing over
$j=1,\ldots,n$, we obtain
\[
\begin{aligned}
b(\bm v_h,\widetilde\rho^n)
={}&
-(\bm\zeta^n-\bm\zeta^0,\bm v_h)
-a_h(\widetilde{\bm\zeta}^n,\bm v_h)
+(\bm w^n-\bm w^0,\bm v_h)
\\
&+\kappa\sum_{j=1}^n\chi_{n-j}
\Bigl[
-(\mathcal R^j,\bm v_h)
+a(\bm w^j,\bm v_h)
\\
&\qquad+
\gamma(\nabla\cdot\bm w^j,
       \nabla\cdot\bm v_h)
-S_h(\bm I_h\bm u^j,\bm v_h)
+
b(\bm v_h,p^j-q_h^j)
\Bigr],
\end{aligned}
\]
where we have used
\[
\kappa\sum_{j=1}^n
\chi_{n-j}\Dt\bm\zeta^j
=
\bm\zeta^n-\bm\zeta^0,
\qquad
\kappa\sum_{j=1}^n
\chi_{n-j}\Dt\bm w^j
=
\bm w^n-\bm w^0.
\]
%
By the discrete inf--sup condition,
the continuity estimates for the bilinear
forms, and Lemma~\ref{lem:interp},
\[
\begin{aligned}
\beta_1\|\widetilde\rho^n\|
\le C\Bigg[
&\|\bm\zeta^n\|+\|\bm\zeta^0\|
+\|\bm w^n\|+\|\bm w^0\|
\\
&+\kappa\sum_{j=1}^n\chi_{n-j}
\Bigl(
M_h\nrm{\bm\zeta^j}
+\|\bm\zeta^j\|
+\|\mathcal R^j\|
+\vartheta_h(t_j)
\Bigr)
\Bigg].
\end{aligned}
\]
%
For brevity, set
\[
E_n=
\|\bm\zeta^0\|
+h^{r+1}
+\tau^{2-\alpha}
+\max_{1\le j\le n}\vartheta_h(t_j).
\]
By Theorem~\ref{thm:fderr} and its proof,
\[
\max_{1\le j\le n}\|\bm\zeta^j\|
+
\left(
\kappa\sum_{j=1}^n
\chi_{n-j}\nrm{\bm\zeta^j}^2
\right)^{1/2}
\le C_{T,\alpha}E_n.
\]
Consequently, the Cauchy--Schwarz inequality
and Lemma~\ref{lem:weights} give
\[
\kappa\sum_{j=1}^n
\chi_{n-j}
\left(
M_h\nrm{\bm\zeta^j}
+\|\bm\zeta^j\|
\right)
\le C_{T,\alpha}(1+M_h)E_n.
\]
%
Moreover, Lemmas~\ref{lem:cons}
and~\ref{lem:weights} yield
\[
\kappa\sum_{j=1}^n
\chi_{n-j}\|\mathcal R^j\|
\le C_{T,\alpha}\tau^{2-\alpha}.
\]
The interpolation estimates imply
$
\|\bm w^n\|+\|\bm w^0\|
\le Ch^{r+1},
$
while
\[
\kappa\sum_{j=1}^n
\chi_{n-j}\vartheta_h(t_j)
\le C_{T,\alpha}
\max_{1\le j\le n}\vartheta_h(t_j).
\]
Combining these bounds gives
\[
\|\widetilde\rho^n\|
\le C_{T,\alpha,\beta_1}(1+M_h)E_n.
\]
%
Finally,
$
P_h^j-p^j
=(P_h^j-q_h^j)+(q_h^j-p^j),
$
and hence
\[
\left\|
\kappa\sum_{j=1}^n
\chi_{n-j}(P_h^j-p^j)
\right\|
\le
\|\widetilde\rho^n\|
+
\kappa\sum_{j=1}^n
\chi_{n-j}\|q_h^j-p^j\|.
\]
By the pressure approximation estimate
and Lemma~\ref{lem:weights}, we have
\[
\kappa\sum_{j=1}^n
\chi_{n-j}\|q_h^j-p^j\|
\le C_{T,\alpha}
\max_{1\le j\le n}\vartheta_h(t_j),
\]
which completes the proof.
\end{proof}



%---------------------------- qqq




\begin{proposition}[Conditional pointwise pressure error]
\label{prop:pointwise-pressure}
Assume the hypotheses of Theorem~\ref{thm:fderr}.
Then, for $1\le n\le N$,
\[
\begin{aligned}
\beta_1\|P_h^n-p^n\|
\le {}&
C_F\|\Dt\bm\zeta^n\|
+M_h\nrm{\bm\zeta^n}
\\
&+C_{T,\alpha}
\left(
\|\bm\zeta^0\|
+h^{r+1}
+\tau^{2-\alpha}
+\max_{1\le j\le n}\vartheta_h(t_j)
\right),
\end{aligned}
\]
where
$
\bm\zeta^n=\bm U_h^n-\bm I_h\bm u^n.
$
The constant $C_{T,\alpha}$ may depend on
$\|\bm b\|_{L^\infty}$ and fixed upper bounds
for $\nu$, $\gamma$, and $\mu$, but is independent
of $h$, $\tau$, and negative powers of $\nu$.
\end{proposition}

%------------------------------


\begin{proof}
Let $q_h^n\in Q_h$ be an approximation of $p^n$.
Rearranging the error equation \eqref{eq:fd-err}
and adding $b(\bm v_h,p^n-q_h^n)$ gives
\[
\begin{aligned}
b(\bm v_h,P_h^n-q_h^n)
={}&
-(\Dt\bm\zeta^n,\bm v_h)
-a_h(\bm\zeta^n,\bm v_h)
\\
&+(\Dt\bm w^n-\mathcal R^n,\bm v_h)
+a(\bm w^n,\bm v_h)
\\
&+\gamma(\nabla\cdot\bm w^n,
         \nabla\cdot\bm v_h)
-S_h(\bm I_h\bm u^n,\bm v_h)
\\
&+b(\bm v_h,p^n-q_h^n),
\end{aligned}
\]
for every $\bm v_h\in\bm V_h$, where
$\bm w^n=\bm u^n-\bm I_h\bm u^n$.
%
By the Friedrichs inequality and the
continuity estimate for $a_h$,
\[
\begin{aligned}
|(\Dt\bm\zeta^n,\bm v_h)|
&\le C_F\|\Dt\bm\zeta^n\|
\|\nabla\bm v_h\|,\\
|a_h(\bm\zeta^n,\bm v_h)|
&\le
\left(
M_h\nrm{\bm\zeta^n}
+\|\bm b\|_{L^\infty}\|\bm\zeta^n\|
\right)
\|\nabla\bm v_h\|.
\end{aligned}
\]
The remaining terms are bounded using
Lemma~\ref{lem:interp} and the continuity
estimates for the bilinear forms. Thus,
\[
\begin{aligned}
|b(\bm v_h,P_h^n-q_h^n)|
\le\Bigl[
&C_F\|\Dt\bm\zeta^n\|
+M_h\nrm{\bm\zeta^n}
\\
&+\|\bm b\|_{L^\infty}\|\bm\zeta^n\|
+C_F\|\Dt\bm w^n\|
\\
&+C_F\|\mathcal R^n\|
+C\vartheta_h(t_n)
\Bigr]\|\nabla\bm v_h\|.
\end{aligned}
\]
Since $P_h^n-q_h^n\in Q_h$, the discrete
inf--sup condition yields
\begin{equation}\label{eq:fd-pressure-intermediate}
\begin{aligned}
\beta_1\|P_h^n-q_h^n\|
\le {}&
C_F\|\Dt\bm\zeta^n\|
+M_h\nrm{\bm\zeta^n}
\\
&+\|\bm b\|_{L^\infty}\|\bm\zeta^n\|
+C_F\|\Dt\bm w^n\|
\\
&+C_F\|\mathcal R^n\|
+C\vartheta_h(t_n).
\end{aligned}
\end{equation}
%
By the interpolation estimate and
Lemma~\ref{lem:cons},
\[
\|\Dt\bm w^n\|
\le C_{T,\alpha}h^{r+1},
\qquad
\|\mathcal R^n\|
\le C_\alpha\tau^{2-\alpha}.
\]
Furthermore, Theorem~\ref{thm:fderr}
and its proof imply
\[
\|\bm\zeta^n\|
\le C_{T,\alpha}
\left(
\|\bm\zeta^0\|
+h^{r+1}
+\tau^{2-\alpha}
+\max_{1\le j\le n}\vartheta_h(t_j)
\right).
\]
%
Substituting these estimates into
\eqref{eq:fd-pressure-intermediate}, we obtain
\[
\begin{aligned}
\beta_1\|P_h^n-q_h^n\|
\le {}&
C_F\|\Dt\bm\zeta^n\|
+M_h\nrm{\bm\zeta^n}
\\
&+C_{T,\alpha}
\left(
\|\bm\zeta^0\|
+h^{r+1}
+\tau^{2-\alpha}
+\max_{1\le j\le n}\vartheta_h(t_j)
\right).
\end{aligned}
\]
%
Finally, the triangle inequality gives
\[
\beta_1\|P_h^n-p^n\|
\le
\beta_1\|P_h^n-q_h^n\|
+\beta_1\|q_h^n-p^n\|.
\]
Using the pressure approximation estimate
\[
\|q_h^n-p^n\|
\le C\vartheta_h(t_n),
\]
and absorbing the resulting term into
the constant proves the assertion.
\end{proof}


%-------------------------------

\begin{remark}\label{rem:final}
Proposition~\ref{prop:fdpressure} provides an estimate
for the fractionally integrated pressure error without
requiring additional pointwise bounds. In contrast,
the pointwise pressure estimate in
Proposition~\ref{prop:pointwise-pressure} is conditional,
since Theorem~\ref{thm:fderr} controls the weighted
energy error
\[
\kappa\sum_{j=1}^n
\chi_{n-j}\nrm{\bm\zeta^j}^2,
\]
but does not provide pointwise bounds for
$\nrm{\bm\zeta^n}$ or $\|\Dt\bm\zeta^n\|$.
If these quantities satisfy
\[
\|\Dt\bm\zeta^n\|
+M_h\nrm{\bm\zeta^n}
\le C_{T,\alpha}
\left(
\|\bm\zeta^0\|
+h^{r+1}
+\tau^{2-\alpha}
+\max_{1\le j\le n}\vartheta_h(t_j)
\right),
\]
then Proposition~\ref{prop:pointwise-pressure}
yields
\[
\|P_h^n-p^n\|
\le C_{T,\alpha,\beta_1}
\left(
\|\bm\zeta^0\|
+h^{r+1}
+\tau^{2-\alpha}
+\max_{1\le j\le n}\vartheta_h(t_j)
\right).
\]
In particular, if
$\bm U_h^0=\bm I_h\bm u^0$,
the resulting convergence rate is
\[
\|P_h^n-p^n\|
=
O\left(
c_hh^r+\gamma^{-1/2}h^{s+1}
+h^{r+1}+\tau^{2-\alpha}
\right).
\]
The numerical experiments presented below examine
the pointwise pressure convergence behaviour
and compare the observed rates with this
conditional estimate.
\end{remark}





%=====================================



\section{Numerical studies}  
\label{sec:numerics}

We present two numerical examples to illustrate the theoretical results and assess the 
performance of the proposed stabilized finite element methods. The first example 
considers a smooth analytical solution, for which we examine spatial and temporal 
convergence separately. The second example considers a solution with weak temporal 
regularity near the initial time and investigates how this reduced regularity affects 
temporal convergence.

All computations are performed on triangular meshes obtained by partitioning
$\Omega=(-1,1)^2$ into an $N\times N$ array of congruent squares and dividing
each square along a diagonal. The resulting mesh consists of $2N^2$ triangles
with mesh size $h=2/N$. For the spatial discretization, we use the inf-sup
stable Taylor--Hood $\mathbb{P}_2/\mathbb{P}_1$ finite element pair.

Grad-div stabilization with parameter $\gamma_{\mathrm{gd}}=1$ is employed in
all computations. We consider the continuous interior penalty (CIP) and local
projection stabilization (LPS) methods, with stabilization parameters
$
\gamma_{\mathrm{cip}}=0.01,
\,
\tau_K=0.01h,
$
respectively. Unless otherwise stated, we take
\[
T=1,\qquad
\nu=10^{-2},\qquad
\boldsymbol b=(1,1)^\top.
\]
% and impose homogeneous Dirichlet boundary conditions for the velocity.

We report the $L^2(\Omega)$ velocity error
$
\|\boldsymbol u_h-\boldsymbol u\|,
$
the velocity-gradient (stress) error
$
\|\nabla(\boldsymbol u_h-\boldsymbol u)\|,
$
the discrete energy-norm error, and the $L^2(\Omega)$ pressure error
$
\|p_h-p\|.
$
All integrals are evaluated using a seven-point Gaussian quadrature rule on
each triangle.



\subsection{Example with a smooth solution}
\label{sec:example1}

Consider the Oseen problem with $\Omega$, $\nu$, $\boldsymbol{b}$, and $T$ as specified above. The right-hand side $\boldsymbol{f}$ is chosen such that
\begin{equation}
\boldsymbol{u}(x,y,t)=t^{2}
\begin{pmatrix}
\sin(\pi x)\cos(\pi y)\\[2pt]
-\cos(\pi x)\sin(\pi y)
\end{pmatrix},
\qquad
p(x,y,t)=t^{2}\sin(\pi x)\sin(\pi y),
\end{equation}
is the exact solution of \eqref{eq:pde}. The solution is smooth in both space and time and satisfies $\boldsymbol{u}(\cdot,0)=\boldsymbol{0}$, making this example suitable for examining the spatial and temporal convergence of the proposed methods.

We first investigate the spatial convergence. We take
\[
N=8,16,32,64, \qquad M=N^2,
\qquad \alpha\in\{0.4,0.6,0.8\},
\]
and use the uniform L1 scheme for the Caputo derivative. With this choice, the temporal error $O(\tau^{2-\alpha})$ is of higher order than the $O(h^2)$ spatial error established in Theorem~\ref{thm:fderr}, allowing the spatial convergence to be observed. The results for the CIP and LPS methods are reported in Tables~\ref{tab:ex1-space-cip} and \ref{tab:ex1-space-lps}, respectively.

The velocity errors exhibit an observed convergence rate close to three for both stabilizations, which is higher than the second-order rate guaranteed by Theorem~\ref{thm:fderr}. The stress, energy-norm, and pressure errors show approximately second-order convergence. The results are consistent across the three values of $\alpha$, indicating that the spatial convergence is essentially unaffected by the fractional order.

We next examine the temporal convergence using the uniform L1
scheme. A sufficiently fine spatial discretization is used so
that the spatial error does not influence the observed temporal
rates. The results for CIP and LPS are presented in
Tables~\ref{tab:ex1-time-cip} and \ref{tab:ex1-time-lps}.
For all three values of $\alpha$, the velocity, stress,
energy-norm, and pressure errors approach the expected rate
$O(\tau^{2-\alpha})$. Thus, the observed temporal orders
approach $1.6$, $1.4$, and $1.2$ for $\alpha=0.4$, $0.6$,
and $0.8$, respectively. The CIP and LPS methods exhibit
very similar errors and convergence behaviour.

Although the pointwise pressure estimate in
Theorem~\ref{thm:fd-pressure-pointwise} contains an additional
first-step contribution, the numerical results exhibit the
optimal temporal rate $2-\alpha$. This behaviour is consistent with the first-step estimate
\[
\|\bm\zeta^1-\bm\zeta^0\|
\le C\tau^\alpha
\bigl(\mathcal E_h+\tau^{2-\alpha}\bigr),
\]
which would control the additional contribution appearing
in the pointwise pressure error bound.

\begin{table}[htbp]
\centering
\caption{Spatial $L^2$ errors and convergence orders for the L1/CIP
method with Taylor--Hood $P_2/P_1$ elements.}
\label{tab:ex1-space-cip}
\setlength{\tabcolsep}{6pt}
\renewcommand{\arraystretch}{1.15}
\begin{tabular}{llcc|cc|cc}
\toprule
\multirow{2}{*}{Quantity} & \multirow{2}{*}{$N_x\times N_y$}
& \multicolumn{2}{c|}{$\alpha=0.4$}
& \multicolumn{2}{c|}{$\alpha=0.6$}
& \multicolumn{2}{c}{$\alpha=0.8$} \\
\cmidrule(lr){3-4}
\cmidrule(lr){5-6}
\cmidrule(l){7-8}
& & $L^2$ error & Order
  & $L^2$ error & Order
  & $L^2$ error & Order \\
\midrule

\multirow{5}{*}{Velocity}
& $4\times4$   & $2.00\times10^{-1}$ & --   & $1.97\times10^{-1}$ & --   & $1.95\times10^{-1}$ & --   \\
& $8\times8$   & $2.49\times10^{-2}$ & 3.01 & $2.48\times10^{-2}$ & 2.99 & $2.50\times10^{-2}$ & 2.96 \\
& $16\times16$ & $3.37\times10^{-3}$ & 2.89 & $3.37\times10^{-3}$ & 2.88 & $3.46\times10^{-3}$ & 2.85 \\
& $32\times32$ & $3.86\times10^{-4}$ & 3.13 & $3.87\times10^{-4}$ & 3.12 & $4.22\times10^{-4}$ & 3.04 \\
& \textit{Theoret. order}
& \multicolumn{2}{c|}{3}
& \multicolumn{2}{c|}{3}
& \multicolumn{2}{c}{3} \\

\addlinespace
\multirow{5}{*}{Stress}
& $4\times4$   & $2.29\times10^{0}$  & --   & $2.26\times10^{0}$  & --   & $2.23\times10^{0}$  & --   \\
& $8\times8$   & $7.34\times10^{-1}$ & 1.64 & $7.32\times10^{-1}$ & 1.63 & $7.30\times10^{-1}$ & 1.61 \\
& $16\times16$ & $2.12\times10^{-1}$ & 1.79 & $2.11\times10^{-1}$ & 1.79 & $2.11\times10^{-1}$ & 1.79 \\
& $32\times32$ & $4.88\times10^{-2}$ & 2.12 & $4.88\times10^{-2}$ & 2.12 & $4.88\times10^{-2}$ & 2.12 \\
& \textit{Theoret. order}
& \multicolumn{2}{c|}{2}
& \multicolumn{2}{c|}{2}
& \multicolumn{2}{c}{2} \\

\addlinespace
\multirow{5}{*}{Energy norm}
& $4\times4$   & $7.26\times10^{-1}$ & --   & $7.21\times10^{-1}$ & --   & $7.17\times10^{-1}$ & --   \\
& $8\times8$   & $1.81\times10^{-1}$ & 2.01 & $1.80\times10^{-1}$ & 2.00 & $1.80\times10^{-1}$ & 2.00 \\
& $16\times16$ & $3.64\times10^{-2}$ & 2.31 & $3.63\times10^{-2}$ & 2.31 & $3.62\times10^{-2}$ & 2.31 \\
& $32\times32$ & $6.17\times10^{-3}$ & 2.56 & $6.17\times10^{-3}$ & 2.56 & $6.17\times10^{-3}$ & 2.56 \\
& \textit{Theoret. order}
& \multicolumn{2}{c|}{2}
& \multicolumn{2}{c|}{2}
& \multicolumn{2}{c}{2} \\

\addlinespace
\multirow{5}{*}{Pressure}
& $4\times4$   & $3.30\times10^{-1}$ & --   & $3.30\times10^{-1}$ & --   & $3.30\times10^{-1}$ & --   \\
& $8\times8$   & $5.92\times10^{-2}$ & 2.48 & $5.92\times10^{-2}$ & 2.48 & $5.92\times10^{-2}$ & 2.48 \\
& $16\times16$ & $1.33\times10^{-2}$ & 2.15 & $1.33\times10^{-2}$ & 2.15 & $1.33\times10^{-2}$ & 2.15 \\
& $32\times32$ & $3.24\times10^{-3}$ & 2.04 & $3.24\times10^{-3}$ & 2.04 & $3.24\times10^{-3}$ & 2.04 \\
& \textit{Theoret. order}
& \multicolumn{2}{c|}{2}
& \multicolumn{2}{c|}{2}
& \multicolumn{2}{c}{2} \\

\bottomrule
\end{tabular}
\end{table}
% =====================================================
% SPATIAL CONVERGENCE: L1/LPS
% =====================================================
\begin{table}[htbp]
\centering
\caption{Spatial $L^2$ errors and convergence orders for the L1/LPS
method with Taylor--Hood $P_2/P_1$ elements.}
\label{tab:ex1-space-lps}
\setlength{\tabcolsep}{6pt}
\renewcommand{\arraystretch}{1.15}
\begin{tabular}{llcc|cc|cc}
\toprule
\multirow{2}{*}{Quantity} & \multirow{2}{*}{$N_x\times N_y$}
& \multicolumn{2}{c|}{$\alpha=0.4$}
& \multicolumn{2}{c|}{$\alpha=0.6$}
& \multicolumn{2}{c}{$\alpha=0.8$} \\
\cmidrule(lr){3-4}
\cmidrule(lr){5-6}
\cmidrule(l){7-8}
& & $L^2$ error & Order
  & $L^2$ error & Order
  & $L^2$ error & Order \\
\midrule

\multirow{5}{*}{Velocity}
& $4\times4$   & $2.29\times10^{-1}$ & --   & $2.24\times10^{-1}$ & --   & $2.21\times10^{-1}$ & --   \\
& $8\times8$   & $3.62\times10^{-2}$ & 2.66 & $3.60\times10^{-2}$ & 2.64 & $3.60\times10^{-2}$ & 2.62 \\
& $16\times16$ & $5.51\times10^{-3}$ & 2.72 & $5.50\times10^{-3}$ & 2.71 & $5.55\times10^{-3}$ & 2.70 \\
& $32\times32$ & $5.78\times10^{-4}$ & 3.25 & $5.78\times10^{-4}$ & 3.25 & $6.02\times10^{-4}$ & 3.20 \\
& \textit{Theoret. order}
& \multicolumn{2}{c|}{3}
& \multicolumn{2}{c|}{3}
& \multicolumn{2}{c}{3} \\

\addlinespace
\multirow{5}{*}{Stress}
& $4\times4$   & $2.68\times10^{0}$  & --   & $2.64\times10^{0}$  & --   & $2.60\times10^{0}$  & --   \\
& $8\times8$   & $1.04\times10^{0}$  & 1.36 & $1.04\times10^{0}$  & 1.34 & $1.03\times10^{0}$  & 1.33 \\
& $16\times16$ & $3.40\times10^{-1}$ & 1.62 & $3.39\times10^{-1}$ & 1.62 & $3.38\times10^{-1}$ & 1.61 \\
& $32\times32$ & $7.25\times10^{-2}$ & 2.23 & $7.24\times10^{-2}$ & 2.23 & $7.23\times10^{-2}$ & 2.22 \\
& \textit{Theoret. order}
& \multicolumn{2}{c|}{2}
& \multicolumn{2}{c|}{2}
& \multicolumn{2}{c}{2} \\

\addlinespace
\multirow{5}{*}{Energy norm}
& $4\times4$   & $7.58\times10^{-1}$ & --   & $7.52\times10^{-1}$ & --   & $7.48\times10^{-1}$ & --   \\
& $8\times8$   & $2.14\times10^{-1}$ & 1.82 & $2.14\times10^{-1}$ & 1.82 & $2.13\times10^{-1}$ & 1.81 \\
& $16\times16$ & $4.68\times10^{-2}$ & 2.20 & $4.67\times10^{-2}$ & 2.19 & $4.65\times10^{-2}$ & 2.19 \\
& $32\times32$ & $7.79\times10^{-3}$ & 2.59 & $7.78\times10^{-3}$ & 2.58 & $7.77\times10^{-3}$ & 2.58 \\
& \textit{Theoret. order}
& \multicolumn{2}{c|}{2}
& \multicolumn{2}{c|}{2}
& \multicolumn{2}{c}{2} \\

\addlinespace
\multirow{5}{*}{Pressure}
& $4\times4$   & $3.31\times10^{-1}$ & --   & $3.31\times10^{-1}$ & --   & $3.31\times10^{-1}$ & --   \\
& $8\times8$   & $5.92\times10^{-2}$ & 2.48 & $5.92\times10^{-2}$ & 2.48 & $5.92\times10^{-2}$ & 2.48 \\
& $16\times16$ & $1.33\times10^{-2}$ & 2.15 & $1.33\times10^{-2}$ & 2.15 & $1.33\times10^{-2}$ & 2.15 \\
& $32\times32$ & $3.24\times10^{-3}$ & 2.04 & $3.24\times10^{-3}$ & 2.04 & $3.24\times10^{-3}$ & 2.04 \\
& \textit{Theoret. order}
& \multicolumn{2}{c|}{2}
& \multicolumn{2}{c|}{2}
& \multicolumn{2}{c}{2} \\

\bottomrule
\end{tabular}
\end{table}

% =====================================================
% TEMPORAL CONVERGENCE: L1/CIP, EXAMPLE 4.1
% =====================================================
\begin{table}[htbp]
\centering
\caption{Temporal $L^2$ errors and convergence orders for the L1/CIP
method with Taylor--Hood $P_2/P_1$ elements.}
\label{tab:ex1-time-cip}
\setlength{\tabcolsep}{6pt}
\renewcommand{\arraystretch}{1.15}
\begin{tabular}{llcc|cc|cc}
\toprule
\multirow{2}{*}{Quantity} & \multirow{2}{*}{$M$}
& \multicolumn{2}{c|}{$\alpha=0.4$}
& \multicolumn{2}{c|}{$\alpha=0.6$}
& \multicolumn{2}{c}{$\alpha=0.8$} \\
\cmidrule(lr){3-4}
\cmidrule(lr){5-6}
\cmidrule(l){7-8}
& & $L^2$ error & Order
  & $L^2$ error & Order
  & $L^2$ error & Order \\
\midrule

\multirow{5}{*}{Velocity}
& 16  & $3.89\times10^{-3}$ & --   & $1.06\times10^{-2}$ & --   & $2.55\times10^{-2}$ & --   \\
& 32  & $1.33\times10^{-3}$ & 1.55 & $4.08\times10^{-3}$ & 1.37 & $1.12\times10^{-2}$ & 1.19 \\
& 64  & $4.47\times10^{-4}$ & 1.57 & $1.56\times10^{-3}$ & 1.39 & $4.86\times10^{-3}$ & 1.20 \\
& 128 & $1.49\times10^{-4}$ & 1.59 & $5.87\times10^{-4}$ & 1.41 & $2.08\times10^{-3}$ & 1.22 \\
& \textit{Theoret. order}
& \multicolumn{2}{c|}{1.6}
& \multicolumn{2}{c|}{1.4}
& \multicolumn{2}{c}{1.2} \\

\addlinespace
\multirow{5}{*}{Stress}
& 16  & $1.73\times10^{-2}$ & --   & $4.70\times10^{-2}$ & --   & $1.13\times10^{-1}$ & --   \\
& 32  & $5.89\times10^{-3}$ & 1.55 & $1.81\times10^{-2}$ & 1.37 & $4.97\times10^{-2}$ & 1.19 \\
& 64  & $1.99\times10^{-3}$ & 1.57 & $6.92\times10^{-3}$ & 1.39 & $2.16\times10^{-2}$ & 1.20 \\
& 128 & $6.62\times10^{-4}$ & 1.59 & $2.61\times10^{-3}$ & 1.41 & $9.24\times10^{-3}$ & 1.22 \\
& \textit{Theoret. order}
& \multicolumn{2}{c|}{1.6}
& \multicolumn{2}{c|}{1.4}
& \multicolumn{2}{c}{1.2} \\

\addlinespace
\multirow{5}{*}{Energy norm}
& 16  & $1.73\times10^{-3}$ & --   & $4.70\times10^{-3}$ & --   & $1.13\times10^{-2}$ & --   \\
& 32  & $5.89\times10^{-4}$ & 1.55 & $1.81\times10^{-3}$ & 1.37 & $4.97\times10^{-3}$ & 1.19 \\
& 64  & $1.99\times10^{-4}$ & 1.57 & $6.92\times10^{-4}$ & 1.39 & $2.16\times10^{-3}$ & 1.20 \\
& 128 & $6.62\times10^{-5}$ & 1.59 & $2.61\times10^{-4}$ & 1.41 & $9.24\times10^{-4}$ & 1.22 \\
& \textit{Theoret. order}
& \multicolumn{2}{c|}{1.6}
& \multicolumn{2}{c|}{1.4}
& \multicolumn{2}{c}{1.2} \\

\addlinespace
\multirow{5}{*}{Pressure}
& 16  & $1.09\times10^{-8}$  & --   & $2.92\times10^{-8}$ & --   & $7.10\times10^{-8}$ & --   \\
& 32  & $3.69\times10^{-9}$  & 1.57 & $1.11\times10^{-8}$ & 1.39 & $3.09\times10^{-8}$ & 1.20 \\
& 64  & $1.24\times10^{-9}$  & 1.58 & $4.23\times10^{-9}$ & 1.40 & $1.33\times10^{-8}$ & 1.21 \\
& 128 & $4.10\times10^{-10}$ & 1.59 & $1.59\times10^{-9}$ & 1.41 & $5.69\times10^{-9}$ & 1.23 \\
& \textit{Theoret. order}
& \multicolumn{2}{c|}{1.6}
& \multicolumn{2}{c|}{1.4}
& \multicolumn{2}{c}{1.2} \\

\bottomrule
\end{tabular}
\end{table}
% =====================================================
% TEMPORAL CONVERGENCE: L1/LPS, EXAMPLE 4.1
% =====================================================
\begin{table}[htbp]
\centering
\caption{Temporal $L^2$ errors and convergence orders for the L1/LPS
method with Taylor--Hood $P_2/P_1$ elements.}
\label{tab:ex1-time-lps}
\setlength{\tabcolsep}{6pt}
\renewcommand{\arraystretch}{1.15}
\begin{tabular}{llcc|cc|cc}
\toprule
\multirow{2}{*}{Quantity} & \multirow{2}{*}{$M$}
& \multicolumn{2}{c|}{$\alpha=0.4$}
& \multicolumn{2}{c|}{$\alpha=0.6$}
& \multicolumn{2}{c}{$\alpha=0.8$} \\
\cmidrule(lr){3-4}
\cmidrule(lr){5-6}
\cmidrule(l){7-8}
& & $L^2$ error & Order
  & $L^2$ error & Order
  & $L^2$ error & Order \\
\midrule

\multirow{5}{*}{Velocity}
& 16  & $3.89\times10^{-3}$ & --   & $1.06\times10^{-2}$ & --   & $2.55\times10^{-2}$ & --   \\
& 32  & $1.33\times10^{-3}$ & 1.55 & $4.08\times10^{-3}$ & 1.37 & $1.12\times10^{-2}$ & 1.19 \\
& 64  & $4.47\times10^{-4}$ & 1.57 & $1.56\times10^{-3}$ & 1.39 & $4.86\times10^{-3}$ & 1.20 \\
& 128 & $1.49\times10^{-4}$ & 1.59 & $5.87\times10^{-4}$ & 1.41 & $2.08\times10^{-3}$ & 1.22 \\
& \textit{Theoret. order}
& \multicolumn{2}{c|}{1.6}
& \multicolumn{2}{c|}{1.4}
& \multicolumn{2}{c}{1.2} \\

\addlinespace
\multirow{5}{*}{Stress}
& 16  & $1.73\times10^{-2}$ & --   & $4.70\times10^{-2}$ & --   & $1.13\times10^{-1}$ & --   \\
& 32  & $5.89\times10^{-3}$ & 1.55 & $1.81\times10^{-2}$ & 1.37 & $4.97\times10^{-2}$ & 1.19 \\
& 64  & $1.99\times10^{-3}$ & 1.57 & $6.92\times10^{-3}$ & 1.39 & $2.16\times10^{-2}$ & 1.20 \\
& 128 & $6.62\times10^{-4}$ & 1.59 & $2.61\times10^{-3}$ & 1.41 & $9.24\times10^{-3}$ & 1.22 \\
& \textit{Theoret. order}
& \multicolumn{2}{c|}{1.6}
& \multicolumn{2}{c|}{1.4}
& \multicolumn{2}{c}{1.2} \\

\addlinespace
\multirow{5}{*}{Energy norm}
& 16  & $1.73\times10^{-3}$ & --   & $4.70\times10^{-3}$ & --   & $1.13\times10^{-2}$ & --   \\
& 32  & $5.89\times10^{-4}$ & 1.55 & $1.81\times10^{-3}$ & 1.37 & $4.97\times10^{-3}$ & 1.19 \\
& 64  & $1.99\times10^{-4}$ & 1.57 & $6.92\times10^{-4}$ & 1.39 & $2.16\times10^{-3}$ & 1.20 \\
& 128 & $6.62\times10^{-5}$ & 1.59 & $2.61\times10^{-4}$ & 1.41 & $9.24\times10^{-4}$ & 1.22 \\
& \textit{Theoret. order}
& \multicolumn{2}{c|}{1.6}
& \multicolumn{2}{c|}{1.4}
& \multicolumn{2}{c}{1.2} \\

\addlinespace
\multirow{5}{*}{Pressure}
& 16  & $1.12\times10^{-8}$  & --   & $3.01\times10^{-8}$ & --   & $7.31\times10^{-8}$ & --   \\
& 32  & $3.80\times10^{-9}$  & 1.57 & $1.15\times10^{-8}$ & 1.39 & $3.18\times10^{-8}$ & 1.20 \\
& 64  & $1.27\times10^{-9}$  & 1.58 & $4.35\times10^{-9}$ & 1.40 & $1.37\times10^{-8}$ & 1.21 \\
& 128 & $4.22\times10^{-10}$ & 1.59 & $1.63\times10^{-9}$ & 1.41 & $5.85\times10^{-9}$ & 1.23 \\
& \textit{Theoret. order}
& \multicolumn{2}{c|}{1.6}
& \multicolumn{2}{c|}{1.4}
& \multicolumn{2}{c}{1.2} \\

\bottomrule
\end{tabular}
\end{table}


% =====================================================
% EXAMPLE 4.2: WEAK REGULARITY
% =====================================================

\subsection{Example with weak temporal regularity}
\label{sec:example2}

The solution considered in Example~4.1 is smooth in time, which is a favourable setting for the uniform L1 scheme. We now consider a solution with reduced temporal regularity by adding a fractional power term to the velocity. With $\Omega$, $\nu$, $\boldsymbol{b}$, and $T$ unchanged, the prescribed solution is
\begin{equation}
\boldsymbol{u}(x,y,t)=\bigl(t^{2}+t^{\alpha}\bigr)
\begin{pmatrix}
\sin(\pi x)\cos(\pi y)\\[2pt]
-\cos(\pi x)\sin(\pi y)
\end{pmatrix},
\qquad
p(x,y,t)=t^{2}\sin(\pi x)\sin(\pi y).
\end{equation}
The term $t^\alpha$ is only H\"older-$\alpha$ regular at $t=0$ and introduces an initial singularity in the time derivatives of the velocity. Consequently, the temporal regularity assumptions used to obtain the $O(\tau^{2-\alpha})$ convergence rate for the uniform L1 scheme are no longer satisfied. 
%Notice, however, that the pressure remains smooth in time.

To account for the reduced temporal regularity, we also consider the graded temporal mesh
\[
t_n=T\left(\frac{n}{M}\right)^r,\qquad 0\leq n\leq M,
\]
where $r\geq1$ is the grading parameter. The choice $r=1$ corresponds to the uniform mesh, while
\[
r=\frac{2-\alpha}{\alpha}
\]
is used to compensate for the initial singularity and recover the optimal temporal convergence rate $O(M^{-(2-\alpha)})$.

We first examine the spatial convergence using
$N=4,8,16,32$ with $M=N^2$ and  $r=\frac{2-\alpha}{\alpha}$.
The results for CIP and LPS are reported in Tables~\ref{tab:ex2-space-cip} and \ref{tab:ex2-space-lps}, respectively. As in the smooth example, the velocity $L^2$-error exhibits an observed convergence rate close to three, while the stress, energy-norm, and pressure errors converge with approximately second order. Thus, the reduced temporal regularity of the velocity does not affect the observed spatial convergence rates when the graded L1 scheme is used.

We next investigate the temporal convergence. Table~\ref{tab:ex2-temporal-cip} presents the results obtained with the uniform L1 scheme. In contrast to the smooth case, the velocity, stress, and energy-norm errors no longer exhibit the $O(\tau^{2-\alpha})$ convergence rate. The observed orders vary with $\alpha$ and with the level of refinement, and a clear asymptotic convergence rate cannot be identified from the reported results. Nevertheless, the deterioration compared with the smooth example is evident and reflects the reduced temporal regularity introduced by the term $t^\alpha$. The pressure error behaves differently: its observed rates remain close to $2-\alpha$. The pressure error behaves differently: its observed rates remain close to $2-\alpha$, which is consistent with the fact that the prescribed pressure remains smooth in time and is therefore not affected by the initial singularity present in the velocity.

Finally, Table~\ref{tab:ex2-temporal-graded-cip} presents the temporal errors obtained with the graded L1 scheme using
$
r=\frac{2-\alpha}{\alpha}.
$
The grading concentrates the time levels near $t=0$, where the velocity has reduced temporal regularity. The numerical results show that the convergence rates of the velocity, stress, and energy-norm errors approach the optimal rate
\[
O(M^{-(2-\alpha)}).
\]
In particular, the observed orders approach $1.6$, $1.4$, and $1.2$ for $\alpha=0.4$, $0.6$, and $0.8$, respectively. The pressure error exhibits the same temporal rates. These results demonstrate that the graded L1 scheme effectively compensates for the initial singularity and restores the optimal temporal convergence behaviour observed in the smooth case.



% =====================================================
% SPATIAL CONVERGENCE: GRADED L1/CIP
% =====================================================
\begin{table}[htbp]
\centering
\caption{Spatial $L^2$ errors and convergence orders for the graded L1/CIP method.}
\label{tab:ex2-space-cip}
\setlength{\tabcolsep}{4pt}
\renewcommand{\arraystretch}{1.1}
\small
\begin{tabular}{llcc|cc|cc}
\toprule
Quantity & $N_x\times N_y$
& \multicolumn{2}{c|}{$\alpha=0.4$}
& \multicolumn{2}{c|}{$\alpha=0.6$}
& \multicolumn{2}{c}{$\alpha=0.8$} \\
\cmidrule(lr){3-4}
\cmidrule(lr){5-6}
\cmidrule(l){7-8}
& & $L^2$ error & Order
& $L^2$ error & Order
& $L^2$ error & Order \\
\midrule

Velocity & $4\times4$ & $4.29\times10^{-1}$ & -- & $4.26\times10^{-1}$ & -- & $4.21\times10^{-1}$ & -- \\
& $8\times8$ & $5.63\times10^{-2}$ & 2.93 & $5.62\times10^{-2}$ & 2.92 & $5.60\times10^{-2}$ & 2.91 \\
& $16\times16$ & $9.06\times10^{-3}$ & 2.64 & $9.05\times10^{-3}$ & 2.63 & $9.05\times10^{-3}$ & 2.63 \\
& $32\times32$ & $1.41\times10^{-3}$ & 2.68 & $1.41\times10^{-3}$ & 2.68 & $1.42\times10^{-3}$ & 2.68 \\

% \addlinespace
% Velocity (max over $t$) & $4\times4$ & $4.29\times10^{-1}$ & -- & $4.26\times10^{-1}$ & -- & $4.21\times10^{-1}$ & -- \\
% & $8\times8$ & $5.63\times10^{-2}$ & 2.93 & $5.62\times10^{-2}$ & 2.92 & $5.60\times10^{-2}$ & 2.91 \\
% & $16\times16$ & $9.06\times10^{-3}$ & 2.64 & $9.05\times10^{-3}$ & 2.63 & $9.05\times10^{-3}$ & 2.63 \\
% & $32\times32$ & $1.41\times10^{-3}$ & 2.68 & $1.41\times10^{-3}$ & 2.68 & $1.42\times10^{-3}$ & 2.68 \\

\addlinespace
Stress & $4\times4$ & $4.94\times10^{0}$ & -- & $4.90\times10^{0}$ & -- & $4.86\times10^{0}$ & -- \\
& $8\times8$ & $1.65\times10^{0}$ & 1.58 & $1.65\times10^{0}$ & 1.57 & $1.64\times10^{0}$ & 1.56 \\
& $16\times16$ & $5.65\times10^{-1}$ & 1.55 & $5.64\times10^{-1}$ & 1.55 & $5.63\times10^{-1}$ & 1.54 \\
& $32\times32$ & $1.79\times10^{-1}$ & 1.66 & $1.79\times10^{-1}$ & 1.66 & $1.79\times10^{-1}$ & 1.66 \\

% \addlinespace
% Stress (max over $t$) & $4\times4$ & $4.94\times10^{0}$ & -- & $4.90\times10^{0}$ & -- & $4.86\times10^{0}$ & -- \\
% & $8\times8$ & $1.65\times10^{0}$ & 1.58 & $1.65\times10^{0}$ & 1.57 & $1.64\times10^{0}$ & 1.56 \\
% & $16\times16$ & $5.65\times10^{-1}$ & 1.55 & $5.64\times10^{-1}$ & 1.55 & $5.63\times10^{-1}$ & 1.54 \\
% & $32\times32$ & $1.79\times10^{-1}$ & 1.66 & $1.79\times10^{-1}$ & 1.66 & $1.79\times10^{-1}$ & 1.66 \\

\addlinespace
Energy norm & $4\times4$ & $1.37\times10^{0}$ & -- & $1.36\times10^{0}$ & -- & $1.36\times10^{0}$ & -- \\
& $8\times8$ & $3.73\times10^{-1}$ & 1.88 & $3.72\times10^{-1}$ & 1.87 & $3.72\times10^{-1}$ & 1.87 \\
& $16\times16$ & $8.33\times10^{-2}$ & 2.16 & $8.32\times10^{-2}$ & 2.16 & $8.31\times10^{-2}$ & 2.16 \\
& $32\times32$ & $1.69\times10^{-2}$ & 2.30 & $1.69\times10^{-2}$ & 2.30 & $1.69\times10^{-2}$ & 2.30 \\

% \addlinespace
% Energy norm (max over $t$) & $4\times4$ & $1.37\times10^{0}$ & -- & $1.36\times10^{0}$ & -- & $1.36\times10^{0}$ & -- \\
% & $8\times8$ & $3.73\times10^{-1}$ & 1.88 & $3.72\times10^{-1}$ & 1.87 & $3.72\times10^{-1}$ & 1.87 \\
% & $16\times16$ & $8.33\times10^{-2}$ & 2.16 & $8.32\times10^{-2}$ & 2.16 & $8.31\times10^{-2}$ & 2.16 \\
% & $32\times32$ & $1.69\times10^{-2}$ & 2.30 & $1.69\times10^{-2}$ & 2.30 & $1.69\times10^{-2}$ & 2.30 \\

\addlinespace
Pressure & $4\times4$ & $3.31\times10^{-1}$ & -- & $3.31\times10^{-1}$ & -- & $3.31\times10^{-1}$ & -- \\
& $8\times8$ & $5.93\times10^{-2}$ & 2.48 & $5.93\times10^{-2}$ & 2.48 & $5.93\times10^{-2}$ & 2.48 \\
& $16\times16$ & $1.33\times10^{-2}$ & 2.15 & $1.33\times10^{-2}$ & 2.15 & $1.33\times10^{-2}$ & 2.15 \\
& $32\times32$ & $3.24\times10^{-3}$ & 2.04 & $3.24\times10^{-3}$ & 2.04 & $3.24\times10^{-3}$ & 2.04 \\

% \addlinespace
% Pressure (max over $t$) & $4\times4$ & $3.31\times10^{-1}$ & -- & $3.31\times10^{-1}$ & -- & $3.31\times10^{-1}$ & -- \\
% & $8\times8$ & $5.93\times10^{-2}$ & 2.48 & $5.93\times10^{-2}$ & 2.48 & $5.93\times10^{-2}$ & 2.48 \\
% & $16\times16$ & $1.33\times10^{-2}$ & 2.15 & $1.33\times10^{-2}$ & 2.15 & $1.33\times10^{-2}$ & 2.15 \\
% & $32\times32$ & $3.24\times10^{-3}$ & 2.04 & $3.24\times10^{-3}$ & 2.04 & $3.24\times10^{-3}$ & 2.04 \\

\bottomrule
\end{tabular}
\end{table}


% =====================================================
% SPATIAL CONVERGENCE: GRADED L1/LPS
% =====================================================
\begin{table}[htbp]
\centering
\caption{Spatial $L^2$ errors and convergence orders for the graded L1/LPS method with Taylor--Hood $P_2/P_1$ elements.}
\label{tab:ex2-space-lps}
\setlength{\tabcolsep}{4pt}
\renewcommand{\arraystretch}{1.1}
\small
\begin{tabular}{llcc|cc|cc}
\toprule
\multirow{2}{*}{Quantity} & \multirow{2}{*}{$N_x\times N_y$}
& \multicolumn{2}{c|}{$\alpha=0.4$}
& \multicolumn{2}{c|}{$\alpha=0.6$}
& \multicolumn{2}{c}{$\alpha=0.8$} \\
\cmidrule(lr){3-4}
\cmidrule(lr){5-6}
\cmidrule(l){7-8}
& & $L^2$ error & Order
& $L^2$ error & Order
& $L^2$ error & Order \\
\midrule

\multirow{5}{*}{Velocity}
& $4\times4$   & $4.29\times10^{-1}$ & --
               & $4.26\times10^{-1}$ & --
               & $4.21\times10^{-1}$ & -- \\
& $8\times8$   & $5.63\times10^{-2}$ & 2.93
               & $5.62\times10^{-2}$ & 2.92
               & $5.60\times10^{-2}$ & 2.91 \\
& $16\times16$ & $9.06\times10^{-3}$ & 2.64
               & $9.05\times10^{-3}$ & 2.63
               & $9.05\times10^{-3}$ & 2.63 \\
& $32\times32$ & $1.41\times10^{-3}$ & 2.68
               & $1.41\times10^{-3}$ & 2.68
               & $1.42\times10^{-3}$ & 2.68 \\
& \textit{Theoret. order}
& \multicolumn{2}{c|}{3}
& \multicolumn{2}{c|}{3}
& \multicolumn{2}{c}{3} \\

% \addlinespace
% \multirow{5}{*}{Velocity (max over $t$)}
% & $4\times4$   & $4.29\times10^{-1}$ & --   & 0.11
%                & $4.26\times10^{-1}$ & --   & 0.02
%                & $4.21\times10^{-1}$ & --   & 0.18 \\
% & $8\times8$   & $5.63\times10^{-2}$ & 2.93 & 1.46
%                & $5.62\times10^{-2}$ & 2.92 & 1.10
%                & $5.60\times10^{-2}$ & 2.91 & 1.41 \\
% & $16\times16$ & $9.06\times10^{-3}$ & 2.64 & 71.34
%                & $9.05\times10^{-3}$ & 2.63 & 82.95
%                & $9.05\times10^{-3}$ & 2.63 & 80.58 \\
% & $32\times32$ & $1.41\times10^{-3}$ & 2.68 & 14966.48
%                & $1.41\times10^{-3}$ & 2.68 & 10181.49
%                & $1.42\times10^{-3}$ & 2.68 & 9770.31 \\
% & \textit{Theoret. order}
% & \multicolumn{3}{c|}{3}
% & \multicolumn{3}{c|}{3}
% & \multicolumn{3}{c}{3} \\

\addlinespace
\multirow{5}{*}{Stress}
& $4\times4$   & $4.94\times10^{0}$  & --
               & $4.90\times10^{0}$  & --
               & $4.86\times10^{0}$  & -- \\
& $8\times8$   & $1.65\times10^{0}$  & 1.58
               & $1.65\times10^{0}$  & 1.57
               & $1.64\times10^{0}$  & 1.56 \\
& $16\times16$ & $5.65\times10^{-1}$ & 1.55
               & $5.64\times10^{-1}$ & 1.55
               & $5.63\times10^{-1}$ & 1.54 \\
& $32\times32$ & $1.79\times10^{-1}$ & 1.66
               & $1.79\times10^{-1}$ & 1.66
               & $1.79\times10^{-1}$ & 1.66 \\
& \textit{Theoret. order}
& \multicolumn{2}{c|}{2}
& \multicolumn{2}{c|}{2}
& \multicolumn{2}{c}{2} \\

% \addlinespace
% \multirow{5}{*}{Stress (max over $t$)}
% & $4\times4$   & $4.94\times10^{0}$  & --   & 0.11
%                & $4.90\times10^{0}$  & --   & 0.02
%                & $4.86\times10^{0}$  & --   & 0.18 \\
% & $8\times8$   & $1.65\times10^{0}$  & 1.58 & 1.46
%                & $1.65\times10^{0}$  & 1.57 & 1.10
%                & $1.64\times10^{0}$  & 1.56 & 1.41 \\
% & $16\times16$ & $5.65\times10^{-1}$ & 1.55 & 71.34
%                & $5.64\times10^{-1}$ & 1.55 & 82.95
%                & $5.63\times10^{-1}$ & 1.54 & 80.58 \\
% & $32\times32$ & $1.79\times10^{-1}$ & 1.66 & 14966.48
%                & $1.79\times10^{-1}$ & 1.66 & 10181.49
%                & $1.79\times10^{-1}$ & 1.66 & 9770.31 \\
% & \textit{Theoret. order}
% & \multicolumn{3}{c|}{2}
% & \multicolumn{3}{c|}{2}
% & \multicolumn{3}{c}{2} \\

\addlinespace
\multirow{5}{*}{Energy norm}
& $4\times4$   & $1.37\times10^{0}$  & --
               & $1.36\times10^{0}$  & --
               & $1.36\times10^{0}$  & -- \\
& $8\times8$   & $3.73\times10^{-1}$ & 1.88
               & $3.72\times10^{-1}$ & 1.87
               & $3.72\times10^{-1}$ & 1.87 \\
& $16\times16$ & $8.33\times10^{-2}$ & 2.16
               & $8.32\times10^{-2}$ & 2.16
               & $8.31\times10^{-2}$ & 2.16 \\
& $32\times32$ & $1.69\times10^{-2}$ & 2.30
               & $1.69\times10^{-2}$ & 2.30
               & $1.69\times10^{-2}$ & 2.30 \\
& \textit{Theoret. order}
& \multicolumn{2}{c|}{2}
& \multicolumn{2}{c|}{2}
& \multicolumn{2}{c}{2} \\

% \addlinespace
% \multirow{5}{*}{Energy norm (max over $t$)}
% & $4\times4$   & $1.37\times10^{0}$  & --   & 0.11
%                & $1.36\times10^{0}$  & --   & 0.02
%                & $1.36\times10^{0}$  & --   & 0.18 \\
% & $8\times8$   & $3.73\times10^{-1}$ & 1.88 & 1.46
%                & $3.72\times10^{-1}$ & 1.87 & 1.10
%                & $3.72\times10^{-1}$ & 1.87 & 1.41 \\
% & $16\times16$ & $8.33\times10^{-2}$ & 2.16 & 71.34
%                & $8.32\times10^{-2}$ & 2.16 & 82.95
%                & $8.31\times10^{-2}$ & 2.16 & 80.58 \\
% & $32\times32$ & $1.69\times10^{-2}$ & 2.30 & 14966.48
%                & $1.69\times10^{-2}$ & 2.30 & 10181.49
%                & $1.69\times10^{-2}$ & 2.30 & 9770.31 \\
% & \textit{Theoret. order}
% & \multicolumn{3}{c|}{2}
% & \multicolumn{3}{c|}{2}
% & \multicolumn{3}{c}{2} \\

\addlinespace
\multirow{5}{*}{Pressure}
& $4\times4$   & $3.31\times10^{-1}$ & --
               & $3.31\times10^{-1}$ & --
               & $3.31\times10^{-1}$ & -- \\
& $8\times8$   & $5.93\times10^{-2}$ & 2.48
               & $5.93\times10^{-2}$ & 2.48
               & $5.93\times10^{-2}$ & 2.48 \\
& $16\times16$ & $1.33\times10^{-2}$ & 2.15
               & $1.33\times10^{-2}$ & 2.15
               & $1.33\times10^{-2}$ & 2.15 \\
& $32\times32$ & $3.24\times10^{-3}$ & 2.04
               & $3.24\times10^{-3}$ & 2.04
               & $3.24\times10^{-3}$ & 2.04 \\
& \textit{Theoret. order}
& \multicolumn{2}{c|}{2}
& \multicolumn{2}{c|}{2}
& \multicolumn{2}{c}{2} \\

% \addlinespace
% \multirow{5}{*}{Pressure (max over $t$)}
% & $4\times4$   & $3.31\times10^{-1}$ & --   & 0.11
%                & $3.31\times10^{-1}$ & --   & 0.02
%                & $3.31\times10^{-1}$ & --   & 0.18 \\
% & $8\times8$   & $5.93\times10^{-2}$ & 2.48 & 1.46
%                & $5.93\times10^{-2}$ & 2.48 & 1.10
%                & $5.93\times10^{-2}$ & 2.48 & 1.41 \\
% & $16\times16$ & $1.33\times10^{-2}$ & 2.15 & 71.34
%                & $1.33\times10^{-2}$ & 2.15 & 82.95
%                & $1.33\times10^{-2}$ & 2.15 & 80.58 \\
% & $32\times32$ & $3.24\times10^{-3}$ & 2.04 & 14966.48
%                & $3.24\times10^{-3}$ & 2.04 & 10181.49
%                & $3.24\times10^{-3}$ & 2.04 & 9770.31 \\
% & \textit{Theoret. order}
% & \multicolumn{3}{c|}{2}
% & \multicolumn{3}{c|}{2}
% & \multicolumn{3}{c}{2} \\

\bottomrule
\end{tabular}
\end{table}
% =====================================================
% TEMPORAL CONVERGENCE: UNIFORM L1/CIP, WEAK REGULARITY
% =====================================================
\begin{table}[htbp]
\centering
\caption{Temporal $L^2$ errors and convergence orders for the graded L1/CIP
method with Taylor--Hood $P_2/P_1$ elements.}
\label{tab:ex2-temporal-graded-cip}
\setlength{\tabcolsep}{2pt}
\renewcommand{\arraystretch}{1.15}
\small
\begin{tabular}{llcc|cc|cc}
\toprule
\multirow{2}{*}{Quantity} & \multirow{2}{*}{$M$}
& \multicolumn{2}{c|}{$\alpha=0.4$}
& \multicolumn{2}{c|}{$\alpha=0.6$}
& \multicolumn{2}{c}{$\alpha=0.8$} \\
\cmidrule(lr){3-4}
\cmidrule(lr){5-6}
\cmidrule(l){7-8}
& & $L^2$ error & Order
  & $L^2$ error & Order
  & $L^2$ error & Order \\
\midrule

% \multirow{6}{*}{Velocity}
% & 4  & $1.67\times10^{-2}$ & --   & $1.53\times10^{-2}$ & --   & $9.22\times10^{-2}$ & --   \\
% & 8  & $1.29\times10^{-2}$ & 0.38 & $5.02\times10^{-3}$ & 1.60 & $3.57\times10^{-2}$ & 1.37 \\
% & 16 & $8.11\times10^{-3}$ & 0.67 & $5.42\times10^{-3}$ & -0.11 & $1.29\times10^{-2}$ & 1.47 \\
% & 32 & $4.61\times10^{-3}$ & 0.82 & $4.03\times10^{-3}$ & 0.43 & $4.25\times10^{-3}$ & 1.60 \\
% & 64 & $2.44\times10^{-3}$ & 0.92 & $2.47\times10^{-3}$ & 0.70 & $1.24\times10^{-3}$ & 1.78 \\
% & \textit{Theoret. order}
% & \multicolumn{2}{c|}{?}
% & \multicolumn{2}{c|}{?}
% & \multicolumn{2}{c}{?} \\

\addlinespace
\multirow{6}{*}{Velocity (max over $t$)}
& 4  & $1.07\times10^{-1}$ & --   & $6.73\times10^{-2}$ & --   & $9.22\times10^{-2}$ & --   \\
& 8  & $9.15\times10^{-2}$ & 0.23 & $5.91\times10^{-2}$ & 0.19 & $3.57\times10^{-2}$ & 1.37 \\
& 16 & $7.35\times10^{-2}$ & 0.32 & $4.46\times10^{-2}$ & 0.40 & $1.66\times10^{-2}$ & 1.10 \\
& 32 & $5.80\times10^{-2}$ & 0.34 & $3.15\times10^{-2}$ & 0.50 & $1.09\times10^{-2}$ & 0.61 \\
& 64 & $4.50\times10^{-2}$ & 0.37 & $2.10\times10^{-2}$ & 0.59 & $6.69\times10^{-3}$ & 0.71 \\
& \textit{Theoret. order}
& \multicolumn{2}{c|}{0.4}
& \multicolumn{2}{c|}{0.6}
& \multicolumn{2}{c}{0.8} \\

% \addlinespace
% \multirow{6}{*}{Stress}
% & 4  & $7.45\times10^{-2}$ & --   & $6.83\times10^{-2}$ & --   & $4.10\times10^{-1}$ & --   \\
% & 8  & $5.72\times10^{-2}$ & 0.38 & $2.25\times10^{-2}$ & 1.60 & $1.59\times10^{-1}$ & 1.37 \\
% & 16 & $3.60\times10^{-2}$ & 0.67 & $2.41\times10^{-2}$ & -0.10 & $5.73\times10^{-2}$ & 1.47 \\
% & 32 & $2.05\times10^{-2}$ & 0.82 & $1.79\times10^{-2}$ & 0.43 & $1.89\times10^{-2}$ & 1.60 \\
% & 64 & $1.08\times10^{-2}$ & 0.92 & $1.10\times10^{-2}$ & 0.70 & $5.52\times10^{-3}$ & 1.78 \\
% & \textit{Theoret. order}
% & \multicolumn{2}{c|}{??}
% & \multicolumn{2}{c|}{??}
% & \multicolumn{2}{c}{??} \\

\addlinespace
\multirow{6}{*}{Stress (max over $t$)}
& 4  & $4.77\times10^{-1}$ & --   & $3.00\times10^{-1}$ & --   & $4.10\times10^{-1}$ & --   \\
& 8  & $4.08\times10^{-1}$ & 0.23 & $2.63\times10^{-1}$ & 0.19 & $1.59\times10^{-1}$ & 1.37 \\
& 16 & $3.28\times10^{-1}$ & 0.31 & $1.99\times10^{-1}$ & 0.40 & $7.42\times10^{-2}$ & 1.10 \\
& 32 & $2.59\times10^{-1}$ & 0.34 & $1.41\times10^{-1}$ & 0.50 & $4.88\times10^{-2}$ & 0.60 \\
& 64 & $2.01\times10^{-1}$ & 0.37 & $9.39\times10^{-2}$ & 0.59 & $2.99\times10^{-2}$ & 0.71 \\
& \textit{Theoret. order}
& \multicolumn{2}{c|}{0.4}
& \multicolumn{2}{c|}{0.6}
& \multicolumn{2}{c}{0.8} \\

% \addlinespace
% \multirow{6}{*}{Energy norm}
% & 4  & $9.99\times10^{-4}$ & --   & $2.14\times10^{-3}$ & --   & $6.02\times10^{-3}$ & --   \\
% & 8  & $6.17\times10^{-4}$ & 0.69 & $8.25\times10^{-4}$ & 1.38 & $2.48\times10^{-3}$ & 1.28 \\
% & 16 & $3.70\times10^{-4}$ & 0.74 & $3.93\times10^{-4}$ & 1.07 & $9.59\times10^{-4}$ & 1.37 \\
% & 32 & $2.07\times10^{-4}$ & 0.84 & $2.16\times10^{-4}$ & 0.86 & $3.72\times10^{-4}$ & 1.36 \\
% & 64 & $1.09\times10^{-4}$ & 0.92 & $1.19\times10^{-4}$ & 0.86 & $1.51\times10^{-4}$ & 1.30 \\
% & \textit{Theoret. order}
% & \multicolumn{2}{c|}{??}
% & \multicolumn{2}{c|}{??}
% & \multicolumn{2}{c}{??} \\

\addlinespace
\multirow{6}{*}{Energy norm (max over $t$)}
& 4  & $7.79\times10^{-3}$ & --   & $5.67\times10^{-3}$ & --   & $6.02\times10^{-3}$ & --   \\
& 8  & $7.72\times10^{-3}$ & 0.01 & $6.14\times10^{-3}$ & -0.12 & $2.48\times10^{-3}$ & 1.28 \\
& 16 & $7.19\times10^{-3}$ & 0.10 & $5.42\times10^{-3}$ & 0.18 & $2.21\times10^{-3}$ & 0.17 \\
& 32 & $6.46\times10^{-3}$ & 0.15 & $4.20\times10^{-3}$ & 0.37 & $1.57\times10^{-3}$ & 0.49 \\
& 64 & $5.55\times10^{-3}$ & 0.22 & $2.93\times10^{-3}$ & 0.52 & $1.03\times10^{-3}$ & 0.61 \\
& \textit{Theoret. order}
& \multicolumn{2}{c|}{0.4}
& \multicolumn{2}{c|}{0.6}
& \multicolumn{2}{c}{0.8} \\

% \addlinespace
% \multirow{6}{*}{Pressure}
% & 4  & $3.79\times10^{-6}$ & --   & $8.06\times10^{-6}$ & --   & $1.53\times10^{-5}$ & --   \\
% & 8  & $1.31\times10^{-6}$ & 1.53 & $3.12\times10^{-6}$ & 1.37 & $6.68\times10^{-6}$ & 1.19 \\
% & 16 & $4.47\times10^{-7}$ & 1.55 & $1.20\times10^{-6}$ & 1.38 & $2.91\times10^{-6}$ & 1.20 \\
% & 32 & $1.51\times10^{-7}$ & 1.57 & $4.55\times10^{-7}$ & 1.39 & $1.26\times10^{-6}$ & 1.21 \\
% & 64 & $5.02\times10^{-8}$ & 1.59 & $1.71\times10^{-7}$ & 1.41 & $5.36\times10^{-7}$ & 1.23 \\
% & \textit{Theoret. order}
% & \multicolumn{2}{c|}{??}
% & \multicolumn{2}{c|}{??}
% & \multicolumn{2}{c}{??} \\

\addlinespace
\multirow{6}{*}{Pressure (max over $t$)}
& 4  & $3.79\times10^{-6}$ & --   & $8.06\times10^{-6}$ & --   & $1.53\times10^{-5}$ & --   \\
& 8  & $1.31\times10^{-6}$ & 1.53 & $3.12\times10^{-6}$ & 1.37 & $6.68\times10^{-6}$ & 1.19 \\
& 16 & $4.47\times10^{-7}$ & 1.55 & $1.20\times10^{-6}$ & 1.38 & $2.91\times10^{-6}$ & 1.20 \\
& 32 & $1.51\times10^{-7}$ & 1.57 & $4.55\times10^{-7}$ & 1.39 & $1.26\times10^{-6}$ & 1.21 \\
& 64 & $5.02\times10^{-8}$ & 1.59 & $1.71\times10^{-7}$ & 1.41 & $5.36\times10^{-7}$ & 1.23 \\
& \textit{Theoret. order}
& \multicolumn{2}{c|}{??}
& \multicolumn{2}{c|}{??}
& \multicolumn{2}{c}{??} \\

\bottomrule
\end{tabular}
\end{table}
% =====================================================
% TEMPORAL CONVERGENCE: GRADED L1/CIP, WEAK REGULARITY
% =====================================================
\begin{table}[htbp]
\centering
\caption{Temporal $L^2$ errors and convergence orders for the graded L1/CIP
method with Taylor--Hood $P_2/P_1$ elements.}
\label{tab:ex2-temporal-graded-cip}
\setlength{\tabcolsep}{2pt}
\renewcommand{\arraystretch}{1.15}
\small
\begin{tabular}{llcc|cc|cc}
\toprule
\multirow{2}{*}{Quantity} & \multirow{2}{*}{$M$}
& \multicolumn{2}{c|}{$\alpha=0.4$}
& \multicolumn{2}{c|}{$\alpha=0.6$}
& \multicolumn{2}{c}{$\alpha=0.8$} \\
\cmidrule(lr){3-4}
\cmidrule(lr){5-6}
\cmidrule(l){7-8}
& & $L^2$ error & Order
  & $L^2$ error & Order
  & $L^2$ error & Order \\
\midrule

% \multirow{6}{*}{Velocity}
% & 4  & $7.55\times10^{-2}$ & --   & $8.50\times10^{-2}$ & --   & $1.24\times10^{-1}$ & --   \\
% & 8  & $3.76\times10^{-2}$ & 1.01 & $3.74\times10^{-2}$ & 1.18 & $5.44\times10^{-2}$ & 1.18 \\
% & 16 & $1.58\times10^{-2}$ & 1.26 & $1.56\times10^{-2}$ & 1.26 & $2.34\times10^{-2}$ & 1.21 \\
% & 32 & $6.02\times10^{-3}$ & 1.39 & $6.29\times10^{-3}$ & 1.31 & $9.98\times10^{-3}$ & 1.23 \\
% & 64 & $2.17\times10^{-3}$ & 1.47 & $2.45\times10^{-3}$ & 1.36 & $4.20\times10^{-3}$ & 1.25 \\
% & \textit{Theoret. order}
% & \multicolumn{2}{c|}{1.6}
% & \multicolumn{2}{c|}{1.4}
% & \multicolumn{2}{c}{1.2} \\

\addlinespace
\multirow{6}{*}{Velocity (max over $t$)}
& 4  & $7.55\times10^{-2}$ & --   & $8.50\times10^{-2}$ & --   & $1.24\times10^{-1}$ & --   \\
& 8  & $3.76\times10^{-2}$ & 1.01 & $3.74\times10^{-2}$ & 1.18 & $5.44\times10^{-2}$ & 1.18 \\
& 16 & $1.58\times10^{-2}$ & 1.26 & $1.56\times10^{-2}$ & 1.26 & $2.34\times10^{-2}$ & 1.21 \\
& 32 & $6.02\times10^{-3}$ & 1.39 & $6.29\times10^{-3}$ & 1.31 & $9.98\times10^{-3}$ & 1.23 \\
& 64 & $2.17\times10^{-3}$ & 1.47 & $2.45\times10^{-3}$ & 1.36 & $4.20\times10^{-3}$ & 1.25 \\
& \textit{Theoret. order}
& \multicolumn{2}{c|}{1.6}
& \multicolumn{2}{c|}{1.4}
& \multicolumn{2}{c}{1.2} \\

% \addlinespace
% \multirow{6}{*}{Stress}
% & 4  & $3.36\times10^{-1}$ & --   & $3.78\times10^{-1}$ & --   & $5.49\times10^{-1}$ & --   \\
% & 8  & $1.67\times10^{-1}$ & 1.01 & $1.66\times10^{-1}$ & 1.18 & $2.42\times10^{-1}$ & 1.18 \\
% & 16 & $7.01\times10^{-2}$ & 1.26 & $6.95\times10^{-2}$ & 1.26 & $1.04\times10^{-1}$ & 1.21 \\
% & 32 & $2.68\times10^{-2}$ & 1.39 & $2.80\times10^{-2}$ & 1.31 & $4.44\times10^{-2}$ & 1.23 \\
% & 64 & $9.65\times10^{-3}$ & 1.47 & $1.09\times10^{-2}$ & 1.36 & $1.87\times10^{-2}$ & 1.25 \\
% & \textit{Theoret. order}
% & \multicolumn{2}{c|}{1.6}
% & \multicolumn{2}{c|}{1.4}
% & \multicolumn{2}{c}{1.2} \\

\addlinespace
\multirow{6}{*}{Stress (max over $t$)}
& 4  & $3.36\times10^{-1}$ & --   & $3.78\times10^{-1}$ & --   & $5.49\times10^{-1}$ & --   \\
& 8  & $1.67\times10^{-1}$ & 1.01 & $1.66\times10^{-1}$ & 1.18 & $2.42\times10^{-1}$ & 1.18 \\
& 16 & $7.01\times10^{-2}$ & 1.26 & $6.95\times10^{-2}$ & 1.26 & $1.04\times10^{-1}$ & 1.21 \\
& 32 & $2.68\times10^{-2}$ & 1.39 & $2.80\times10^{-2}$ & 1.31 & $4.44\times10^{-2}$ & 1.23 \\
& 64 & $9.65\times10^{-3}$ & 1.47 & $1.09\times10^{-2}$ & 1.36 & $1.87\times10^{-2}$ & 1.25 \\
& \textit{Theoret. order}
& \multicolumn{2}{c|}{1.6}
& \multicolumn{2}{c|}{1.4}
& \multicolumn{2}{c}{1.2} \\

% \addlinespace
% \multirow{6}{*}{Energy norm}
% & 4  & $4.95\times10^{-3}$ & --   & $6.08\times10^{-3}$ & --   & $8.38\times10^{-3}$ & --   \\
% & 8  & $2.47\times10^{-3}$ & 1.00 & $2.74\times10^{-3}$ & 1.15 & $3.79\times10^{-3}$ & 1.14 \\
% & 16 & $1.03\times10^{-3}$ & 1.26 & $1.15\times10^{-3}$ & 1.26 & $1.64\times10^{-3}$ & 1.21 \\
% & 32 & $3.91\times10^{-4}$ & 1.40 & $4.59\times10^{-4}$ & 1.32 & $7.05\times10^{-4}$ & 1.22 \\
% & 64 & $1.40\times10^{-4}$ & 1.48 & $1.78\times10^{-4}$ & 1.36 & $3.02\times10^{-4}$ & 1.22 \\
% & \textit{Theoret. order}
% & \multicolumn{2}{c|}{1.6}
% & \multicolumn{2}{c|}{1.4}
% & \multicolumn{2}{c}{1.2} \\

\addlinespace
\multirow{6}{*}{Energy norm (max over $t$)}
& 4  & $4.95\times10^{-3}$ & --   & $6.08\times10^{-3}$ & --   & $8.38\times10^{-3}$ & --   \\
& 8  & $2.47\times10^{-3}$ & 1.00 & $2.76\times10^{-3}$ & 1.14 & $3.79\times10^{-3}$ & 1.14 \\
& 16 & $1.03\times10^{-3}$ & 1.26 & $1.37\times10^{-3}$ & 1.00 & $1.64\times10^{-3}$ & 1.21 \\
& 32 & $3.91\times10^{-4}$ & 1.40 & $6.10\times10^{-4}$ & 1.17 & $7.05\times10^{-4}$ & 1.22 \\
& 64 & $1.40\times10^{-4}$ & 1.48 & $2.50\times10^{-4}$ & 1.28 & $3.45\times10^{-4}$ & 1.03 \\
& \textit{Theoret. order}
& \multicolumn{2}{c|}{1.6}
& \multicolumn{2}{c|}{1.4}
& \multicolumn{2}{c}{1.2} \\

% \addlinespace
% \multirow{6}{*}{Pressure}
% & 4  & $1.63\times10^{-5}$ & --   & $1.92\times10^{-5}$ & --   & $2.24\times10^{-5}$ & --   \\
% & 8  & $7.50\times10^{-6}$ & 1.12 & $8.48\times10^{-6}$ & 1.18 & $1.03\times10^{-5}$ & 1.13 \\
% & 16 & $3.01\times10^{-6}$ & 1.32 & $3.49\times10^{-6}$ & 1.28 & $4.58\times10^{-6}$ & 1.17 \\
% & 32 & $1.12\times10^{-6}$ & 1.43 & $1.38\times10^{-6}$ & 1.33 & $2.01\times10^{-6}$ & 1.19 \\
% & 64 & $3.96\times10^{-7}$ & 1.50 & $5.34\times10^{-7}$ & 1.37 & $8.62\times10^{-7}$ & 1.22 \\
% & \textit{Theoret. order}
% & \multicolumn{2}{c|}{1.6}
% & \multicolumn{2}{c|}{1.4}
% & \multicolumn{2}{c}{1.2} \\

\addlinespace
\multirow{6}{*}{Pressure (max over $t$)}
& 4  & $1.63\times10^{-5}$ & --   & $1.92\times10^{-5}$ & --   & $2.24\times10^{-5}$ & --   \\
& 8  & $7.50\times10^{-6}$ & 1.12 & $8.48\times10^{-6}$ & 1.18 & $1.03\times10^{-5}$ & 1.13 \\
& 16 & $3.01\times10^{-6}$ & 1.32 & $3.49\times10^{-6}$ & 1.28 & $4.58\times10^{-6}$ & 1.17 \\
& 32 & $1.12\times10^{-6}$ & 1.43 & $1.38\times10^{-6}$ & 1.33 & $2.01\times10^{-6}$ & 1.19 \\
& 64 & $3.96\times10^{-7}$ & 1.50 & $5.34\times10^{-7}$ & 1.37 & $8.62\times10^{-7}$ & 1.22 \\
& \textit{Theoret. order}
& \multicolumn{2}{c|}{1.6}
& \multicolumn{2}{c|}{1.4}
& \multicolumn{2}{c}{1.2} \\

\bottomrule
\end{tabular}
\end{table}

% =====================================================
% =====================================================================
\begin{thebibliography}{99}
\bibitem{AKJ} N.~Ahmed, S.~Karaa, A.~Jha, Symmetric stabilized FEM for time-fractional convection--diffusion--reaction equations, \emph{Math.\ Comput.\ Simulation} \textbf{245} (2026), 685--697.
\bibitem{Ali} A.~A.~Alikhanov, A priori estimates for solutions of boundary value problems for fractional-order equations, \emph{Differ.\ Equ.} \textbf{46} (2010), 660--666.
\bibitem{BBF} D.~Boffi, F.~Brezzi, M.~Fortin, \emph{Mixed Finite Element Methods and Applications}, Springer, Berlin, 2013.
{\color{blue}
\bibitem{Baz01} E.~Bazhlekova, \emph{Fractional Evolution Equations in Banach Spaces}, PhD thesis, Eindhoven University of Technology, 2001.
\bibitem{BF07} E.~Burman, M.~A.~Fern\'andez, Continuous interior penalty finite element method for the time-dependent Navier--Stokes equations: space discretization and convergence, \emph{Numer.\ Math.} \textbf{107} (2007), 39--77.
\bibitem{Die10} K.~Diethelm, \emph{The Analysis of Fractional Differential Equations}, Lecture Notes in Math.\ 2004, Springer, Berlin, 2010.
\bibitem{GS03} V.~Girault, L.~R.~Scott, A quasi-local interpolation operator preserving the discrete divergence, \emph{Calcolo} \textbf{40} (2003), 1--19.
\bibitem{LL18} L.~Li, J.-G.~Liu, A generalized definition of Caputo derivatives and its application to fractional ODEs, \emph{SIAM J.\ Math.\ Anal.} \textbf{50} (2018), 2867--2900.
\bibitem{YGD07} H.~Ye, J.~Gao, Y.~Ding, A generalized Gronwall inequality and its application to a fractional differential equation, \emph{J.\ Math.\ Anal.\ Appl.} \textbf{328} (2007), 1075--1081.
}
\bibitem{dFGJN} J.~de~Frutos, B.~Garc\'{\i}a-Archilla, V.~John, J.~Novo, Analysis of the grad-div stabilization for the time-dependent Navier--Stokes equations with inf-sup stable finite elements, \emph{Adv.\ Comput.\ Math.} \textbf{44} (2018), 195--225.
\bibitem{Kal} T.~Kaluza, \"Uber die Koeffizienten reziproker Potenzreihen, \emph{Math.\ Z.} \textbf{28} (1928), 161--170.
\bibitem{LX} Y.~Lin, C.~Xu, Finite difference/spectral approximations for the time-fractional diffusion equation, \emph{J.\ Comput.\ Phys.} \textbf{225} (2007), 1533--1552.
\bibitem{Lin} A.~Linke, On the role of the Helmholtz decomposition in mixed methods for incompressible flows and a new variational crime, \emph{Comput.\ Methods Appl.\ Mech.\ Engrg.} \textbf{268} (2014), 782--800.
\end{thebibliography}

% =====================================================================
\appendix
\section{Proof of the consistency estimate for the L1 scheme}\label{app:cons}
\begin{proof}[Proof of \eqref{lem:cons}]
Write $M_2=\max_{t\in[0,T]}\|\partial_t^2\bm u(t)\|_X$ and $\bm u^k=\bm u(t_k)$.

Since $\bm u\in C^1([0,T];X)$, splitting the integral in \eqref{eq:caputo} at the nodes gives
\[
\CD\bm u(t_n)=\frac1{\Gamma(1-\alpha)}\sum_{k=1}^n\int_{t_{k-1}}^{t_k}(t_n-s)^{-\alpha}\,\bm u'(s)\,ds .
\]
For the L1 approximation, we use
\[
\int_{t_{k-1}}^{t_k}(t_n-s)^{-\alpha}\,ds=\frac{(t_n-t_{k-1})^{1-\alpha}-(t_n-t_k)^{1-\alpha}}{1-\alpha}=\frac{\tau^{1-\alpha}\bigl((n-k+1)^{1-\alpha}-(n-k)^{1-\alpha}\bigr)}{1-\alpha},
\]
together with $\Gamma(2-\alpha)=(1-\alpha)\Gamma(1-\alpha)$, which shows
\[
\tau^{-\alpha}\beta_{n-k}\,(\bm u^k-\bm u^{k-1})=\frac1{\Gamma(1-\alpha)}\,\frac{\bm u^k-\bm u^{k-1}}\tau\int_{t_{k-1}}^{t_k}(t_n-s)^{-\alpha}\,ds .
\]
Summing over $k$ and subtracting the two representations,
\begin{equation}\label{eq:cons-integral}
\mathcal R^n=-\frac1{\Gamma(1-\alpha)}\sum_{k=1}^n\int_{t_{k-1}}^{t_k}\varphi_n(s)\,\bm e_k(s)\,ds,
\end{equation}
with
\[
\varphi_n(s):=(t_n-s)^{-\alpha},\qquad \bm e_k(s):=\bm u'(s)-\frac{\bm u^k-\bm u^{k-1}}\tau .
\]

Since $(\bm u^k-\bm u^{k-1})/\tau=\tau^{-1}\int_{t_{k-1}}^{t_k}\bm u'(\theta)\,d\theta$,
\begin{equation}\label{eq:cons-ek}
\int_{t_{k-1}}^{t_k}\bm e_k(s)\,ds=0,\qquad
\bm e_k(s)=\frac1\tau\int_{t_{k-1}}^{t_k}\bigl(\bm u'(s)-\bm u'(\theta)\bigr)\,d\theta,
\end{equation}
and, because $\|\bm u'(s)-\bm u'(\theta)\|_X\le M_2|s-\theta|\le M_2\tau$ for $s,\theta\in[t_{k-1},t_k]$, \eqref{eq:cons-ek} gives
\begin{equation}\label{eq:cons-ek-bound}
\|\bm e_k(s)\|_X\le M_2\tau .
\end{equation}

On the last interval $k=n$, $\varphi_n$ is singular at $s=t_n$ but integrable, and \eqref{eq:cons-ek-bound} gives
\begin{equation}\label{eq:cons-last}
\Bigl\|\int_{t_{n-1}}^{t_n}\varphi_n(s)\,\bm e_n(s)\,ds\Bigr\|_X\le M_2\tau\int_{t_{n-1}}^{t_n}(t_n-s)^{-\alpha}\,ds=\frac{M_2\,\tau^{2-\alpha}}{1-\alpha}.
\end{equation}

On the earlier intervals $k\le n-1$, $\varphi_n$ is smooth on $[t_{k-1},t_k]$ with $\varphi_n'(s)=\alpha(t_n-s)^{-\alpha-1}>0$, which is increasing, so
\[
\sup_{s\in[t_{k-1},t_k]}|\varphi_n'(s)|=\alpha\bigl((n-k)\tau\bigr)^{-\alpha-1}.
\]
By the first relation in \eqref{eq:cons-ek}, the constant $\varphi_n(t_k)$ can be subtracted:
\[
\int_{t_{k-1}}^{t_k}\varphi_n(s)\,\bm e_k(s)\,ds=\int_{t_{k-1}}^{t_k}\bigl(\varphi_n(s)-\varphi_n(t_k)\bigr)\,\bm e_k(s)\,ds .
\]
With $|\varphi_n(s)-\varphi_n(t_k)|\le\sup|\varphi_n'|\,(t_k-s)$, $\int_{t_{k-1}}^{t_k}(t_k-s)\,ds=\tau^2/2$ and \eqref{eq:cons-ek-bound},
\begin{equation}\label{eq:cons-earlier}
\Bigl\|\int_{t_{k-1}}^{t_k}\varphi_n\,\bm e_k\,ds\Bigr\|_X\le M_2\tau\cdot\alpha\bigl((n-k)\tau\bigr)^{-\alpha-1}\cdot\frac{\tau^2}2=\frac\alpha2\,M_2\,\tau^{2-\alpha}\,(n-k)^{-1-\alpha}.
\end{equation}

Since $\sum_{j=1}^{\infty}j^{-1-\alpha}\le1+\int_1^\infty x^{-1-\alpha}\,dx=1+\frac1\alpha$, summing \eqref{eq:cons-earlier} over $k=1,\dots,n-1$ gives
\begin{equation}\label{eq:cons-sum}
\sum_{k=1}^{n-1}\frac\alpha2\,M_2\tau^{2-\alpha}(n-k)^{-1-\alpha}\le\frac\alpha2\Bigl(1+\frac1\alpha\Bigr)M_2\tau^{2-\alpha}=\frac{1+\alpha}2\,M_2\,\tau^{2-\alpha}.
\end{equation}
Combining \eqref{eq:cons-integral} with the bounds \eqref{eq:cons-last} and \eqref{eq:cons-sum} for its $n$ summands, and multiplying by $1/\Gamma(1-\alpha)$,
\[
\|\mathcal R^n\|_X\le\frac1{\Gamma(1-\alpha)}\Bigl(\frac{1+\alpha}2+\frac1{1-\alpha}\Bigr)M_2\,\tau^{2-\alpha}=C_\alpha M_2\,\tau^{2-\alpha}. \qedhere
\]
\end{proof}


\begin{remark}
The proof uses $\bm u\in C^2([0,T];X)$ up to $t=0$. If $\partial_t^2\bm u$ is unbounded near $t=0$ (initial fractional singularity), the first intervals contribute a larger error and weighted or graded-mesh consistency bounds must replace \eqref{lem:cons}.
\end{remark}


\end{document}